% 自 問答.md 由 tools/md_to_tex.py 於 2026-11-02 轉出；部內順序依歸部表。此後本檔為正本。

\qaentry{2026-10-11}{Kelsen 與維也納學派}

\begin{qaprompt}
Hans Kelsen 不就是維也納學派嗎
\end{qaprompt}

\begin{qaquestion}
Kelsen 與「維也納學派」有什麼關係？本週兩篇在哪裡提到這個學派？
\end{qaquestion}

Vinx 譯本 Kelsen 一篇的篇首註寫明底本：Klecatsky、Marcic、Schambeck 合編《Die Wiener rechtstheoretische Schule. Schriften von Hans Kelsen, Adolf Merkl, Alfred Verdross》（Verlag Österreich, 2010）第 2 冊頁 1485--531。\footcite[22]{kelsen2015}書名所稱的維也納法理論學派，這套文集收的是 Kelsen、Merkl、Verdross 三人的著作。

Schmitt 在 94 頁註 30 引 A. Hensel 刊於 \emph{Archiv des öffentlichen Rechts} 新編第 15 卷（1928）415 頁的一句話：「a monarchist, then, would be someone whose opinion differs from that of the Vienna School」。Schmitt 說這句玩笑點中了要害，形式主義的公法學正可以帶有特定的政治意義。\footcite[94--95]{schmitt2015}同一個註的前一句，引的就是 Kelsen 本週這篇的原刊 81 頁。

Kelsen 在 23 至 25 頁把法的創設排成憲法、法律、命令、裁判、強制執行的階層，每一層對上一層是法的適用，對下一層是法的創設。Schmitt 在 109 至 111 頁註 58 批評的「continuous hierarchy」，指的就是這套講法。

\qaentry{2026-10-12}{Kelsen 與 Gödel：時間重疊、同校但不同學院與學術圈}

\begin{qaprompt}
同時期嗎？同校園？同學術環境？
\end{qaprompt}

\begin{qaquestion}
Hans Kelsen 與 Kurt Gödel 是否同時代、同在維也納大學，以及處於同一學術環境？
\end{qaquestion}

\subsection{精確回答}

\begin{itemize}
\tightlist
\item
  同時期：是。Kelsen 生於 1881 年；Gödel 生於 1906 年。Gödel 1924 年進入維也納大學時，Kelsen 已是維也納法學界的重要教授；Kelsen 1930 年離開維也納，Gödel 仍在維也納大學任職至 1930 年代。
\item
  同校園：是，至少有明確的時間重疊。Kelsen 屬法學與國法學／法理論；Gödel 屬哲學院體系中的數學、邏輯與科學哲學，不是同一學院，也沒有證據顯示兩人是師生。
\item
  同學術環境：可以說共享「維也納」的知識與思想環境，但不能說是同一專業學術圈。Kelsen 在維也納法學院形成維也納法理論學派；Gödel 約 1926 年加入維也納學派（Vienna Circle），由數學家 Hans Hahn 引介。兩個圈子有城市與跨學科思想上的鄰近性，但研究問題、訓練與學術網絡不同。
\end{itemize}

因此最精確的表述是：Kelsen 與 Gödel 是同一時代、曾在同一所大學和同一座城市活動的不同學科人物；他們共享部分維也納學術氣候，但沒有資料支持「同樣受訓」或「同一學派」的說法。

維也納大學官方資料：Kelsen 1911 年在維也納大學取得國法與法哲學 habilitation，1930 年離開維也納；Gödel 1924 年入學，1930 年取得博士並於同年在該校任職，約 1926 年加入維也納學派。

\qaentry{2026-10-12}{Kelsen 的形式法學與 Gödel：同一維也納學術環境，不是同樣訓練}

\begin{qaprompt}
我自己的疑問：Hans Kelsen？形式法學？是不是和哥德爾受同樣訓練？
\end{qaprompt}

\begin{qaquestion}
Kelsen 的「形式法學」是否與 Gödel 的形式化、邏輯訓練有共同來源？兩人是否接受過相同的學術訓練？
\end{qaquestion}

\subsection{先給結論}

不是同樣訓練；比較準確的說法是：兩人共享一部分維也納的學術與思想環境，但分屬不同學科。

\begin{itemize}
\tightlist
\item
  Hans Kelsen：法律學、國法學與法哲學；他的核心工作是把法學建構成研究有效法規範、位階、授權與法創設過程的「純粹法學」。
\item
  Kurt Gödel：先在維也納大學學物理，後轉向數學；他的核心工作是數理邏輯、數學基礎與形式系統的可證明性問題。
\end{itemize}

\subsection{Kelsen 所謂「形式」是什麼}

Kelsen 的形式性不是把法律化約成數學，也不是只研究法條的外觀。它是把法秩序當成一個由規範位階與授權關係構成的結構，追問：

\begin{itemize}
\tightlist
\item
  下位規範是否依據上位規範產生；
\item
  哪一機關依何種程序有權創設規範；
\item
  法律的效力如何沿著憲法、法律、命令、裁判與執行逐層具體化。
\end{itemize}

W05 預習原文正是這個意思：立法與執行不是兩個並列功能，而是法創設過程中上下位階的階段；合法性是下位階與上位階法秩序相符的關係。\footcites[23--24]{kelsen2015}[31--32]{kelsen1929}因此，「形式法學」較準確地理解為結構與效力關係的法學，不是沒有內容的形式主義。

\subsection{和 Gödel 可以怎麼比較}

可以作一個方法上的類比：兩人都重視形式結構、有效性條件與系統內部的推導關係。但不能因此說 Kelsen 是按照 Gödel 式數理邏輯受訓，也不能把 Gödel 的不完備性定理直接套到 Kelsen 的法秩序上。

兩人的交集首先是維也納的時代與學術環境。維也納大學官方資料記載，Kelsen 的領域是法律、憲法與國際法、法理論；Gödel 則於 1924 年進入維也納大學，先學物理，後轉向數學、邏輯與科學哲學。兩人因此是同一大學與同一城市思想環境中的不同世代、不同學科人物，而不是接受同一套專業訓練。

\subsection{與本週文本的關係}

Kelsen 的形式性最後服務於一個憲法制度問題：如果法律效力是分層產生的，就必須有一個獨立機關確認下位規範是否符合上位憲法，並以撤銷違憲行為保障法秩序。這使他的憲法法院理論帶有高度結構性，但憲法法院實際處理的仍是法律規範與權限問題，而不是 Gödel 意義下的形式系統證明。

\qaentry{2026-10-12}{「事實情境涵攝進規則」的文本位置與 Karl Kraus 的維也納脈絡}

\begin{qaprompt}
文本定位：事實情境涵攝進規則

我們卡老師當時在哪？Karl Kraus
\end{qaprompt}

\subsection{文本定位：Schmitt，Vinx 英譯本 83、114、116--117 頁}

你說的「事實情境涵攝進規則」，最直接的文本在 Schmitt 討論 1925 年帝國法院判決的段落：法官把具體案件的事實，涵攝到憲法律或普通法律所規定的構成要件／規範之下。這是普通法院的典型司法活動，德文術語是 Subsumtion（涵攝）。\footcites[15]{schmitt1931}[83]{schmitt2015}

但要和後面的兩種「涵攝」分開：

\begin{itemize}
\tightlist
\item
  事實 → 規則：具體事實涵攝於法律規範，這是典型司法裁判；
\item
  規範 → 規範：比較憲法律與普通法律，決定兩者是否衝突、哪一個有效，Schmitt 說這裡根本不是涵攝（\emph{Es wird überhaupt nicht subsumiert}）；\footcites[43]{schmitt1931}[114]{schmitt2015}
\item
  不確定概念 → 確定內容：例如判斷神學院是否屬於憲法所稱的神學院；若以有拘束力的方式消除憲法內容疑義，Schmitt 認為實質上是立法甚至憲法立法，而非司法。\footcites[44--45]{schmitt1931}[116--117]{schmitt2015}
\end{itemize}

所以你要的那句應整理成：法官把事實情境涵攝進既有規則，是司法；憲法法院若先替規範內容作一般性、權威性的確定，則可能已不是單純司法。

這裡的關鍵可再精確寫成：只有在憲法規範已經提供足夠確定的規則，使具體事實可以涵攝於其下時，法院的工作才是司法性的涵攝。 如果法院必須先決定憲法規範本身究竟包含什麼，才知道事實是否落入規範，那就不再只是「憲法規範允許的涵攝」，而是對憲法內容的權威確定。

因此，憲法規範無法進行「規範對規範」的涵攝：普通法律不是被涵攝進憲法規範，而是與憲法規範相互比較，以確認兩者是否衝突及何者有效。Schmitt 的原文是：「\emph{Es wird überhaupt nicht subsumiert}」〔根本沒有進行涵攝〕。\footcites[43]{schmitt1931}[114]{schmitt2015}

Schmitt 的神學院例正是在分界這兩件事：是否把某一具體學院涵攝為憲法所稱的神學院，表面上像事實涵攝；但若這個判斷其實是在決定憲法概念的內容，法院便是在消除規範疑義，實質上進入立法甚至憲法立法。\footcites[114, 116--117]{schmitt2015}[43--45]{schmitt1931}

\subsection{Karl Kraus 在哪裡}

Karl Kraus（1874--1936）不是 Kelsen 或 Gödel 的同學，也不是同一學術專業的人；他是維也納的作家、記者、諷刺家與《Die Fackel》創辦人。可是他與兩人確實屬於同一時代的維也納知識與文化環境：Kelsen 在維也納法學院形成維也納法理論學派，Gödel 約 1926 年加入維也納學派，Kraus 則透過《Die Fackel》批判新聞語言、政治修辭與社會文化。

因此 Kraus 適合放在「維也納共同文化場域」的位置，不宜放成 Kelsen 的法學老師、Gödel 的邏輯老師，或與他們同一學派。奧地利國家圖書館將 Kraus 定位為 1874--1936 年人物，並記載《Die Fackel》作為其反對煽情報刊與語言敗壞的媒介；這正是他與 Kelsen、Gödel 所共享的時代背景，但不是同一訓練。

\qaentry{2026-10-12}{Kelsen 談的是德國還是奧地利？}

\begin{qaprompt}
Hans Kelsen 講的是德國還是奧地利？？
\end{qaprompt}

\begin{qaquestion}
Kelsen〈Wesen und Entwicklung der Staatsgerichtsbarkeit〉的制度背景與論證對象，主要是德國還是奧地利？這篇文章與 Schmitt〈Der Hüter der Verfassung〉的國別脈絡有何不同？
\end{qaquestion}

Kelsen 不是在做單一國家的制度介紹，而是以奧地利的制度經驗提出一般性的憲法裁判理論。他本人是奧地利法學家，主張由集中於單一機關的憲法法院審查法律，制度上的重要背景是 1920 年奧地利憲法法院模式。文章的理論出發點因此較接近奧地利，而不是德國帝國法院的附隨審查。

但文章的發表場域是德國法學界：Kelsen 是奧地利學者，受邀在德國國家法學者協會於維也納舉行的會議中作共同報告；德文原刊收於 \emph{Veröffentlichungen der Vereinigung der Deutschen Staatsrechtslehrer} 第 5 輯（1929）。他也處理德國與奧地利共同涉及的問題，例如在國際法作為憲法審查基準的討論中並列兩國。\footcites[58]{kelsen2015}[66]{kelsen1929}所以較準確的說法是：Kelsen 以奧地利模式為重要制度背景，向德國法學界提出可一般化的憲法裁判理論。

\subsection{Kelsen：憲法是什麼的原則？}

Kelsen 在狹義上把憲法界定為：決定整個法秩序與國家秩序的最高原則；進一步說，是決定法律如何創設、由哪些機關依何種程序創設法律的規則。

\begin{quote}
德文：„so ergibt sich die Vorstellung eines obersten, die ganze Rechts- und Staatsordnung bestimmenden, für das Wesen der ... Gemeinschaft entscheidenden Prinzips.``

英文：``what results is the idea of a highest principle that determines the whole legal and political order, a principle that is decisive for the nature of the community constituted by that order.''

中譯：由此得到的，是一個最高原則的觀念：它決定整個法律秩序與政治秩序，並決定由該秩序所構成之共同體的本質。
\end{quote}

其後的狹義定義是：

\begin{quote}
德文：„Dies: die Regel für die Erzeugung der die staatliche Ordnung vor allem bildenden Rechtsnormen, die Bestimmung der Organe und des Verfahrens der Gesetzgebung, ist der eigentliche, ursprüngliche und engere Begriff der Verfassung.``

英文：``This rule for the creation of the legal norms that primarily constitute the state order, the determination of the organs and the procedure of legislation, is the actual, original and narrower concept of the constitution.''

中譯：規定構成國家秩序之主要法律規範如何創設，並規定立法機關與立法程序的規則，才是憲法本來、原初而較狹義的概念。\footcites[36]{kelsen1929}[27--28]{kelsen2015}
\end{quote}

\subsection{\texorpdfstring{憲法在法律中的具體化（\emph{Konkretisierung}）}{憲法在法律中的具體化（Konkretisierung）}}

Kelsen 所說的不是「憲法內容被一般法律簡單重複」，而是法規範沿著位階逐層具體化：憲法規定法律創設的框架，法律在此框架內創設命令，命令再由裁判或行政行為適用到個案，最後落實為強制執行。每一層對上位階是法的適用或再生產，對下位階又是法的創設。

\begin{quote}
德文：„Der Weg, den das Recht von der Verfassung bis zu dem die richterliche oder Verwaltungsentscheidung vollziehenden Tatbestand nimmt, ist ein solcher steter Konkretisierungsprozeß.``

英文：``The route that the law travels, on its way from the constitution to the matter of fact that enforces a judicial or administrative decision, is one of continuous concretization.''

中譯：法律從憲法一路走到使司法或行政決定獲得執行的事實狀態，所經歷的是一個持續的具體化過程。
\end{quote}

因此，「憲法在法律中具體化」的意思是：憲法的規範要透過下位階法律取得較具體的內容；但下位階法律仍受憲法拘束，不能把具體化理解成法律可以自由改寫憲法。Kelsen 以此說明法律對憲法的合憲性，以及命令、裁判與行政行為對上位階規範的合法性。\footcites[23--25]{kelsen2015}[31--33]{kelsen1929}

\subsection{\texorpdfstring{\emph{monarchical principle}（\emph{monarchisches Prinzip}；君主原則）}{monarchical principle（monarchisches Prinzip；君主原則）}}

Kelsen 所說的 \emph{monarchical principle} 不是單純描述「有君主的制度」，而是十九世紀君主立憲國法學中的一項政治---法理學說：把君主視為剩餘國家權力的中心，將國會參與立法理解為對君主決定的「同意」（\emph{Zustimmung}），並由君主的簽署與公布確認法律的成立與合憲性。

Kelsen 認為，這套原則不是憲法條文自身的中立說明，而是以君主原則重新解釋實證憲法，因而使君主掌握本應受到控制的合憲性控制功能；部長副署也不能有效補救這個問題。\footcites[25--27]{kelsen2015}[33--35]{kelsen1929}因此，較準確的中譯是君主原則，不是泛稱的「君主制原則」。

這裡的「實證憲法」（\emph{positive constitution}；德文 \emph{positive Verfassung}）是指作為現行有效法規範而存在的憲法：包括憲法所規定的立法機關、立法程序與其他具有憲法位階的規範。它不是「經驗上觀察到的憲法」，也不是「實證研究所得的憲法」。Kelsen 的批評是，君主立憲學說不依實證憲法本身，而以外加的 \emph{monarchical principle} 重新解釋它；因此，必須先區分實證憲法的規範內容與君主原則的政治---法理學說。

\subsection{Kelsen 文中的副署}

Kelsen 提到副署時，並不是在主張「內閣以副署箝制一個虛位總統」。他是在批判立憲君主制的傳統理論：該理論把君主簽署與公布法律視為法律合憲性的保障，但 Kelsen 認為，主張控制功能的恰恰是「本身就需要受控制的機關」。\footcites[26--27]{kelsen2015}[35]{kelsen1929}在這個脈絡下，部長副署只是使君主的行為附有部長責任；若問題出在國會的立法程序，部長副署也不能提供有效的控制。

關鍵句三語對照如下：

\begin{quote}
德文：„Allein die Kontrollfunktion wird gerade von jener Instanz beansprucht, die selbst der Kontrolle bedürfte.``

英文：``The function of control is claimed precisely by the one power that is itself most in need of being controlled.''

中譯：然而，主張控制功能的，恰恰是本身最需要受控制的那個權力機關。
\end{quote}

同一段對副署的說法是：

\begin{quote}
德文：„Zwar wird der Akt des Monarchen durch die ministerielle Gegenzeichnung unter Verantwortung gestellt.``

英文：``Admittedly, the act of the monarch is put under responsibility, by virtue of the requirement of ministerial countersignature.''

中譯：固然，依部長副署的要求，君主的行為被置於責任之下。
\end{quote}

Kelsen 後面承認，部長責任與副署可以作為追究立法程序合憲性的「個人責任」保證，但它仍然不充分，因為即使追究了個人責任，違憲法律本身仍可能繼續有效。\footcites[43--44]{kelsen2015}[60--61]{kelsen1929}因此，Kelsen 的結論是：副署是較弱的責任性保障，真正有效的憲法保障仍是由憲法法院撤銷違憲法律。

相對地，Schmitt 第一章明確以威瑪德國為主要對象，討論德國帝國法院（Reichsgericht）、威瑪憲法及誰應擔任德國憲法的守護者。\footcite[79--90]{schmitt2015}兩篇放在一起讀時，可概括為：Kelsen 從奧地利式集中型憲法法院出發，Schmitt 則針對德國的附隨審查與憲政秩序提出批判。

\qaentry{2026-10-12}{立法與執行：不是本質二分，而是法創設過程的上下階段}

\begin{qaprompt}
執行和立法？什麼沒有本質的差別，只有高低、前後的差別
\end{qaprompt}

\begin{qaquestion}
Kelsen 是否認為立法與執行沒有本質差別？「高低、前後」具體指什麼？
\end{qaquestion}

\subsection{Kelsen 的原文}

這裡的「執行」不是日常語言中單指把既有命令付諸實施，而是 \emph{Vollziehung}：包括立法以下的法適用，尤其命令、行政處分、裁判，最後才到強制執行。Kelsen 明確說：

\begin{quote}
„Gesetzgebung und Vollziehung sind nicht zwei koordinierte Staatsfunktionen, sondern sind nur zwei im Verhältnis von Unter- und Überordnung stehende Stufen des Prozesses der Rechtserzeugung.``

「立法與執行不是兩種並列的國家功能，而只是法創設過程中處於上下位階關係的兩個階段。」
\end{quote}

所以，不是說立法與執行在制度上完全相同，而是說兩者沒有不可跨越的本質差別：

\begin{itemize}
\tightlist
\item
  立法相對位於較高階段，依憲法創設一般規範，法形成的自由較大；
\item
  執行相對位於較低階段，依上位規範創設較具體的命令、裁判或行政處分，受上位規範拘束較強。
\end{itemize}

每一階段都同時具有兩面：對上一階段而言是法的適用，對下一階段而言又是法的創設。因此，立法不是純粹自由創造，執行也不是機械重述；差別是法創設過程中的相對高低與前後位置，以及自由與拘束的比例不同。\footcites[22--24]{kelsen2015}[30--32]{kelsen1929}

\subsection{這和「消極／積極立法者」的關係}

這個階層理論說明為什麼 Kelsen 能把憲法法院稱為消極立法者：法院撤銷一般規範時，具有一般性的法創設效果，但它是在更高位階憲法的拘束下作成裁判；國會則是在憲法框架內有較大形成空間的積極立法者。兩者的差異不是「一個創設法律、另一個完全不創設法律」，而是創設的方向、位階與拘束程度不同。

\qaentry{2026-10-12}{Kelsen：法規範位階、緊急命令與人的保障}

\begin{qaprompt}
法律位階的法規範／緊急命令／行政行為欠缺法律依據／\emph{persönliche Garantien}
\end{qaprompt}

\begin{qaquestion}
Kelsen 如何把不同位階的法規範、緊急命令、欠缺法律依據的行政行為，以及副署制度下的人的保障，放在同一個憲法保障架構中？
\end{qaquestion}

\subsection{不同法律位階的法規範}

較準確的說法是不同法律位階的法規範（\emph{Rechtsnormen verschiedener Stufen}）。Kelsen 把憲法、法律、命令、裁判與執行行為安排在法秩序的不同層級：下位規範由上位規範授權或拘束，並在此基礎上進一步具體化。違憲或違法的判斷，就是檢查某一行為是否符合較高位階的規範。\footcites[23--25, 56]{kelsen2015}[31--33, 65]{kelsen1929}

放回我國法脈絡，《憲法訴訟法》的法規範憲法審查區分法律位階法規範與命令位階法規範，且兩者的違憲審查與失效效果不完全相同（第 49、52、55、59 條；本機《憲法訴訟法》整理）。但這是我國憲法訴訟的程序分類，不能直接等同於 Kelsen 從憲法排列到個別執行行為的完整法階層：前者回答「憲法法庭審查哪一類法規範」，後者回答「法規範如何創設、拘束與具體化」。

\subsection{\texorpdfstring{緊急命令（\emph{Notverordnung}）}{緊急命令（Notverordnung）}}

Kelsen 所說的緊急命令不是一般行政機關執行法律的命令，而是憲法直接授權、可以代替或變更法律的命令。它直接位於憲法之下，法律效力與法律並列，因此也可能直接違憲。

\begin{quote}
德文：„Es sind dies insbesondere die sog. Notverordnungen.``

英文：``The so-called emergency decrees, in particular, belong to this class.''

中譯：尤其包括所謂的緊急命令。
\end{quote}

放回我國法，總統緊急命令的憲法依據、行政院會議決議、立法院追認與不同意時的失效，須依憲法增修條文及憲法第 43 條的制度脈絡分別判斷，不能直接以 Kelsen 的 \emph{Notverordnung} 取代我國規範（本機憲法增修條文整理）。

\subsection{行政行為欠缺法律依據}

行政機關作成行為時，若完全沒有法律依據，或明顯濫引法律，該行為即不只是違反個別法律，而是直接違反憲法上的執行合法性原則。Kelsen 原文說：„sondern er ist gesetzlos und daher als solcher unmittelbar verfassungswidrig.``（「而是欠缺法律依據，因此本身即直接違憲。」）\footcites[41]{kelsen1929}[33]{kelsen2015}

\subsection{\texorpdfstring{\emph{persönliche Garantien} 與副署}{persönliche Garantien 與副署}}

Kelsen 把合法性保障分為預防的／壓制的，以及人的／物的保障（\emph{präventive／repressive; persönliche／sachliche Garantien}）。\emph{persönliche Garantien}（personal guarantees；人的保障）是對作成公權力行為之人的刑事、懲戒或民事責任，例如元首與部長的責任；\emph{sachliche Garantien}（物的保障）則是違法行為本身的無效或撤銷。\footcites[35--36, 43--44]{kelsen2015}[43--44, 51--52]{kelsen1929}

副署放在人的保障的脈絡理解：它使部長對元首行為承擔政治責任。Kelsen 的重點是，人的保障本身不足以消除違法行為；真正有效的憲法保障仍是由獨立機關宣告違憲行為無效或予以撤銷。\footcites[43--44]{kelsen2015}[51--52]{kelsen1929}因此，不能把「行政行為欠缺法律依據」與「不副署」寫成同一命題；前者是行為的法源與合憲性問題，後者是責任與程序問題，兩者的制度連結須依具體憲法規定判斷。

\subsection{\texorpdfstring{\emph{Kassation}：廢棄，不是「撤銷法院」}{Kassation：廢棄，不是「撤銷法院」}}

這裡要區分名詞與功能：

\begin{itemize}
\tightlist
\item
  \emph{Kassationsgericht} 是廢棄法院（亦可譯為破棄法院），不是「撤銷法院」；
\item
  \emph{Kassation} 是法院廢棄某一判決或法律的權能；
\item
  \emph{kassatorisch} 是廢棄性的，描述判決直接除去違法規範的效果。
\end{itemize}

Kelsen 說憲法法院的判決必須具有「廢棄的性質」（\emph{kassatorischen Charakter}）。因此，在本篇脈絡中，\emph{Kassation} 可依對象譯為「廢棄法律」或「撤銷違憲法律」；但法院本身不能譯成「撤銷法院」。放回我國法用語，最高法院通常「廢棄原判決」，行政法院通常「撤銷行政處分」，憲法法庭則依《憲法訴訟法》宣告法規範違憲及失效，三者不要混譯。\footcites[61--64]{kelsen2015}[70--73]{kelsen1929}

\qaentry{2026-10-12}{副署與虛位元首}

\begin{qaprompt}
副署是內閣箝制虛位的總統
\end{qaprompt}

\begin{qaquestion}
副署制度是否就是內閣用來箝制虛位元首的制度？這個理解能否直接套用到我國總統與行政院院長的關係？
\end{qaquestion}

這句話抓到副署制度的古典功能，但還需要補上「政治責任」這一層。副署的作用不只是限制元首，而是使元首公布法律、發布命令時，必須有對國會負責的部長或內閣首長共同承擔政治責任；因此，副署制度同時處理元首的法律行為效力與內閣的責任歸屬。

我國憲法第三十七條規定，總統公布法律、發布命令，須經行政院院長，或行政院院長及有關部會首長副署。這確實使行政院院長具有制約總統行為的制度位置；但憲法增修條文另規定，總統發布行政院院長及依法須經立法院同意任命人員的任免命令，以及解散立法院的命令，無須行政院院長副署。因此，我國不能簡化為「內閣箝制一個虛位總統」。

較準確的說法是：副署在我國保留了內閣制中由行政院院長承擔政治責任、並制約總統部分行為的功能；但總統由人民直接選舉，且享有提名行政院院長等憲法權限，總統並非單純的虛位元首。司法院釋字第 419 號曾以「總統代表國家，統而不治；行政院院長領導政府，治而不統」說明原始憲法的設計；釋字第 627 號則指出修憲後總統直選、行政院院長改由總統任命，中央政府體制已具有不同於原始設計的結構。

\qaentry{2026-10-12}{Kelsen 的刑事責任：人的保障與憲法法院撤銷的區別}

\begin{qaprompt}
刑事責任
\end{qaprompt}

\begin{qaquestion}
刑事責任在 Kelsen 的憲法保障理論中扮演什麼角色？它能否取代憲法法院撤銷違憲規範？
\end{qaquestion}

Kelsen 把刑事責任放在人的保障（\emph{persönliche Garantien}）中：由元首、部長或其他具體行為人，對違反憲法或法律的行為承擔刑事、懲戒或政治責任。副署制度與部長責任即屬此類保障；它把責任歸屬到可以被追究的個人，而不是抽象地歸給國家機關。

但國會是合議機關，通常不適合承擔刑事或民事責任；國會議員的集體決議也不能因此直接轉化為每一名議員的刑事責任。相對地，元首與部長可以在制度上成為責任主體。\footcites[43--44]{kelsen2015}[51--52]{kelsen1929}

刑事責任的限制在於：它主要制裁行為人，卻不必然除去已經成立並正在適用的違憲法律。即使部長被追究責任，違憲法律仍可能留在法秩序中繼續有效。因此，Kelsen 認為人的保障是不充分的；真正有效的憲法保障，是由獨立於立法者的憲法法院直接撤銷違憲行為。\footcites[43--44]{kelsen2015}[51--52]{kelsen1929}

所以要分開：

\begin{itemize}
\tightlist
\item
  刑事責任：追究誰作成或副署了違法、違憲行為，是人的、責任性的保障。
\item
  憲法法院撤銷：除去違憲的一般規範本身，是物的、壓制性的保障。
\end{itemize}

前者處理「誰要負責」，後者處理「違憲規範是否繼續有效」；兩者不能互相替代。

\qaentry{2026-10-12}{Kelsen 的「命令」與我國法規命令違憲審查}

\begin{qaprompt}
大法官的解釋，審查「法規命令」的違憲
\end{qaprompt}

\begin{qaquestion}
Kelsen 為什麼把命令納入憲法裁判的客體？這如何對應到我國大法官／憲法法庭對法規命令的違憲審查？
\end{qaquestion}

\subsection{Kelsen 文本的脈絡}

Kelsen 在列舉憲法法院應審查的客體時，先談法律，接著談命令。\footcites[48--51]{kelsen2015}[57--60]{kelsen1929}他不是把所有行政命令一概當作憲法問題，而是先區分兩類：

\begin{itemize}
\tightlist
\item
  代替或變更法律的緊急命令：直接依據憲法，位於憲法之下、與法律並列；政府若藉此越過國會的法律保留與權限界線，直接涉及憲法問題。
\item
  執行法律的命令：內容以法律為上位規範，首先是命令是否合法；但 Kelsen 認為，法律合憲性與命令合法性的審查具有「內在親緣關係」，在具體憲法安排下也可交由憲法法院審查。
\end{itemize}

所以他的位階鏈是：憲法 → 法律 → 命令。直接依據憲法的命令，以憲法為審查基準；執行法律的命令，以法律為主要審查基準。命令是否同時涉及憲法，要看它是直接執行憲法、代替法律，還是單純執行法律。\footcites[58--65]{kelsen1929}[50--56]{kelsen2015}

\subsection{對應我國法}

舊制大法官解釋時期，命令就是違憲審查的對象；《司法院大法官審理案件法》第 14 條第 1 項曾明定，宣告命令牴觸憲法，以出席大法官過半數同意行之。現行《憲法訴訟法》則使用「命令位階法規範」的分類，法規命令可以成為法規範憲法審查對象；法規命令除不得牴觸憲法，也不得逾越法律授權或牴觸法律。

法規命令與行政規則不能只依名稱區分：法規命令是依法授權、對外發生法律效果的一般抽象規範；行政規則原則上是機關內部規範。但若行政機關以「行政規則」之名實際對外規制人民，仍須依規範的實質內容判斷其法律性質。

\qaentry{2026-10-12}{Kelsen 文中的合議庭與合議機關}

\begin{qaprompt}
合議／不合議
\end{qaprompt}

\begin{qaquestion}
Kelsen 文中的合議庭、合議機關與獨任制，分別在什麼脈絡下出現？
\end{qaquestion}

正式用語是合議制／合議機關與獨任制／獨任機關；但 Kelsen 在本文至少有兩個不同脈絡：一是憲法法院作為法院，由多名法官共同評議；二是國會作為合議機關，涉及責任能否歸屬於個別成員。

在憲法法院的脈絡，Kelsen 說的是 \emph{Gerichtskollegium}（法院合議庭）：憲法法院由多名成員組成，對法律問題共同評議；他主張成員不宜過多，並討論國會、政府、法律專家等不同組成方式。\footcites[47--48, 67]{kelsen2015}[56--57, 76--77]{kelsen1929}原文並說：„die Öffentlichkeit der Beratung des Urteils durch das Gerichtskollegium selbst``，即「法院合議庭對判決的評議是否公開」。

在人的保障的脈絡，Kelsen 說國會是合議機關，通常不適合承擔刑事或民事責任；元首與部長則可以負責。這裡的合議機關不是在討論法院是否採合議庭，而是在討論集體機關的責任歸屬。\footcites[43--44]{kelsen2015}[51--52]{kelsen1929}放回我國法，憲法法庭屬合議機關；若要表達單一首長作成決定，才用獨任制，若要表達未經共同評議，則寫未經合議程序。

\qaentry{2026-10-12}{消極立法者、積極立法者與 \emph{actio popularis}}

\begin{qaprompt}
消極立法者／積極立法者；\emph{action popularis}
\end{qaprompt}

\begin{qaquestion}
Kelsen 為什麼把國會稱為積極立法者、憲法法院稱為消極立法者？\emph{actio popularis} 在這個架構中又是什麼？
\end{qaquestion}

\subsection{Kelsen 的區分}

國會是積極立法者（\emph{positiver Gesetzgeber}），在憲法所劃定的框架內，選擇並創設一般規範的內容；它雖受憲法拘束，仍保有較大的政治形成空間。憲法法院則是消極立法者（\emph{negativer Gesetzgeber}）：它不是重新選擇政策、制定替代法律，而是在憲法所決定的基準下，將違憲法律自法秩序中廢棄。因此 Kelsen 說：

\begin{quote}
„Der negative Gesetzgeber aber, das Verfassungsgericht, ist bei seiner Funktion wesentlich durch die Verfassung bestimmt.``

「消極立法者，也就是憲法法院，其功能在本質上由憲法決定。」
\end{quote}

這個比喻不是說憲法法院與國會具有相同的立法權，而是說：廢棄一般規範的效果，在功能上具有立法性；但其裁判活動主要是適用憲法，拘束程度高於國會的積極形成活動。所以 Kelsen 接著說，憲法法院的功能在這個意義上更接近一般法院，是真正的裁判。\footcites[47]{kelsen2015}[56]{kelsen1929}

這個區分也有一個重要例外：若憲法法院不只廢棄違憲規範，還宣告已被廢止的舊規範恢復效力，法院的功能就會「帶有更強的立法性質」。但這仍不是法院自由創設新政策，而是對規範失效後法秩序如何銜接所作的憲法裁判安排。\footcites[63]{kelsen2015}[73]{kelsen1929}

\subsection{\texorpdfstring{\emph{Actio popularis}：最強的開啟審查方式}{Actio popularis：最強的開啟審查方式}}

\emph{Actio popularis}（拉丁文；德文本文亦作 \emph{actio popularis}；中譯為民眾訴訟）是指：不要求聲請人證明自己的權利受到個別侵害，只要任何人提出聲請，憲法法院即須開啟規範審查程序。它使消極立法者最容易發現並廢棄違憲法律，因而是對憲法保障最強的程序設計。

但 Kelsen 並不推薦採用。他的理由不是這種制度保障力不足，而是可能造成濫訴，使憲法法院負擔過重。\footcites[64--65]{kelsen2015}[74]{kelsen1929}他因此另外討論法院提請、訴訟當事人間接挑戰、國會少數、憲法檢察官及法院依職權審查等較受限制的開啟方式。

\subsection{與我國憲法訴訟的區分}

我國人民聲請憲法審查不是一般性的 \emph{actio popularis}。依《憲法訴訟法》第 59 條，人民須主張憲法上權利遭受不法侵害，並經通常救濟程序、對不利確定終局裁判及其所適用的法規範提出聲請；這是個案權利受侵害所觸發的法規範憲法審查，不是任何人無須利害關係即可挑戰任何法律。

因此要分開兩個問題：聲請程序是否由個案或公益開啟，以及憲法法庭作成違憲判決後，法規範失效的效力是否具有一般性。我國人民聲請欠缺 \emph{actio popularis} 的公益起訴資格，不代表法規範違憲判決的效果必然只限於聲請人；後者仍依《憲法訴訟法》第 51、52 條及判決主文決定。

\qaentry{2026-10-12}{憲法法庭的失效、民法上的無效與法安定性}

\begin{qaprompt}
失效，幾種失效的模式／大法官的無效和民法的無效／法安定性／例外情況：個案，人民聲請解釋
\end{qaprompt}

\begin{qaquestion}
憲法法庭宣告法規範違憲的「無效／失效」，與民法上法律行為的「無效」有何不同？有哪些失效模式？人民以個案聲請憲法審查時，效果是否只限於該個案？
\end{qaquestion}

\subsection{憲法法庭的用語：宣告違憲不必然等於立即失效}

我國《憲法訴訟法》先區分「宣告法規範違憲」與「宣告違憲且應失效」。法規範被宣告違憲後，是否失效、何時失效與是否溯及，須看判決主文；不能把每一個「違憲宣告」都直接翻成民法上的「自始無效」（《憲法訴訟法》第 51、52 條）。

\subsection{法規範失效的主要模式}

\begin{itemize}
\tightlist
\item
  判決生效日起失效（通常所稱立即失效）：判決主文沒有另定時點時，法規範自判決生效日起失效。
\item
  溯及失效：判決主文指定法規範在較早時點失效；判決前已確定的裁判，仍須依《憲法訴訟法》第 53 條區分其是否得依法定程序救濟。
\item
  定期失效：法規範在判決後仍暫時維持效力，留給立法者修法期間；法律位階法規範期限不得逾二年，命令位階法規範不得逾一年。
\item
  違憲但不宣告失效：判決只確認規範與憲法不合，未使其立即退出法秩序；這時須依判決主文與理由判斷後續機關的處理義務。
\end{itemize}

因此，「無效／失效」至少要問三件事：對象是什麼、效力從何時起、既有裁判與個案如何處理。這也是 Kelsen 所說法安定性問題的制度化表現：違憲審查不能只問規範有無瑕疵，還要安排撤銷的時間與範圍。

\subsection{大法官／憲法法庭的失效，與民法上的無效}

民法上的無效，主要是法律行為在成立或內容上違反民法所定要件的效果：違反強制或禁止規定、違背公序良俗、欠缺法定方式時，原則上無效（民法第 71 至 73 條）；顯失公平則是法院得依聲請撤銷或減輕給付（民法第 74 條），不是當然無效。

憲法法庭處理的則是法規範或裁判是否違反憲法，以及宣告後的失效、廢棄與救濟效果。它不是把民法第 71 條的「法律行為無效」直接套到法律或命令上；尤其法規範在判決前可能已有效適用，判決也可能為定期失效或只對將來失效。故「民法上的無效」與「憲法法庭宣告法規範失效」是不同層次的效力制度。

\subsection{法安定性與人民個案聲請}

法安定性要求保護既有法律關係、確定裁判與人民對現行法秩序的信賴，所以《憲法訴訟法》允許定期失效，並規定立即失效原則上不影響判決前已作成的確定裁判；溯及失效時，才依法律或判決意旨開啟救濟（第 52 至 54 條）。

人民依《憲法訴訟法》第 59 條聲請，須先經通常救濟程序，針對不利確定終局裁判及其所適用的法規範提出憲法審查。聲請具有個案起因，但效果不必然只限於聲請人：憲法法庭若認所適用的法規範違憲，仍依第 51 至 54 條決定法規範的失效方式；若認確定終局裁判違憲，則依第 62 條宣告裁判違憲並廢棄、發回管轄法院。個案性主要表現在聲請門檻與原因案件救濟，不等於憲法法庭的法規範失效永遠只有個案效力。

\subsection{通案與個案：消除什麼問題}

通案是法規範憲法審查：處理規範本身對不特定多數案件的普遍違憲問題，必要時宣告法規範失效。個案是裁判憲法審查：處理人民在特定案件中因確定終局裁判所受的憲法權利侵害，透過宣告裁判違憲、廢棄並發回，消除該個案的憲法侵害。

因此，「消除個案的不正義」可以作為課堂上的直觀說法，但正式法學表述應是救濟個案中的憲法權利侵害；它不是泛指法院重新判斷哪一方較不公平，也不排除同一聲請同時導致所適用法規範發生一般性的失效效果。

\qaentry{2026-10-12}{第五次修憲失效與 \emph{reformieren}}

\begin{qaprompt}
第五次憲法增修條文違憲、無效，回到上一次的釋字 599；\emph{reformation}
\end{qaprompt}

\begin{qaquestion}
第五次修憲被宣告違憲後，是否回到前一次的增修條文？這和 Kelsen 所說的 \emph{reformieren} 是同一種效果嗎？
\end{qaquestion}

這裡的解釋編號是釋字第 499 號。釋字第 499 號宣告 1999 年第五次修憲中相關修正條文自解釋公布日起失其效力，並明文指出：1997 年 7 月 21 日修正公布的原增修條文繼續適用。所以這不是抽象地說「違憲修憲當然全部回到上一次」，而是本號解釋對特定修正條文指定失效時點與替代適用的規範狀態。

Kelsen 所說的 \emph{reformieren} 是變更／改作合法行為：撤銷機關可以自己把瑕疵行為改成合法內容，也可以只撤銷後交由原機關重作。\footcites[43]{kelsen2015}[51]{kelsen1929}釋字第 499 號不是憲法法庭行使這種 \emph{reformieren}，而是宣告第五次修憲相關條文失效，並讓先前的增修條文繼續適用；其性質較接近失效後的規範回復／繼續適用效果，不是法院替修憲機關重新制定內容。

\subsection{釋字第499號與 Kelsen 的理論漏洞}

這裡的「Kelsen 漏洞」不是說 Kelsen 忘了修憲，而是指其純粹法學的位階模型會遇到一個界限問題：如果修憲法律依形式上有效的修憲程序制定，它究竟由哪一個更高位階的規範審查？ 若憲法本身沒有明文不可修改的核心或更高規範，修憲法律一旦有效成立，就很難在同一國內法秩序中再以「違憲」否定它。

釋字第499號採取的是臺灣憲法實務的回應：憲法修改雖屬修憲機關權限，仍受正當修憲程序與憲法基本原則拘束；本號並因第五次修憲程序重大瑕疵，宣告相關修正條文自解釋公布日起失效，並使先前增修條文繼續適用。這不是 Kelsen 本人提出的解法，而是以憲法解釋建立對修憲權的實質與程序界限，補足形式位階模型的不足。

但必須保留一個區分：Kelsen 在本週文章中並非完全沒有處理更高規範問題；他另討論國際法優位時，承認修憲法律也可能因牴觸條約或一般國際法而有瑕疵。那是以國際法優位為前提的另一條效力鏈，不等於臺灣釋字第499號所建立的國內修憲界限。

\qaentry{2026-10-12}{Kelsen 的形式法學與 Schmitt 的批評}

\begin{qaprompt}
Kelsen 的形式法學？Carl Schmitt 非常有意見
\end{qaprompt}

\begin{qaquestion}
Kelsen 所謂的形式法學，形式究竟在哪裡？Schmitt 為什麼強烈反對？
\end{qaquestion}

\subsection{Kelsen 的「形式」}

Kelsen 的形式性，不是指法律沒有內容，而是先問：規範由誰、依何種程序、在何種位階關係中創設，並由什麼較高規範授權與拘束。憲法、法律、命令、裁判與行政行為形成一個法創設的階層過程；每一階段對上位階是法的適用，對下位階又是法的創設。\footcites[23--25]{kelsen2015}[31--33]{kelsen1929}

因此，Kelsen 的憲法法院理論不是只審查立法程序是否完成。若憲法對法律內容有明確規定，憲法法院也可審查內容是否抵觸憲法；但審查基準必須來自實證憲法，而不能由法官自行引入未實證化的正義、自然法或合目的性標準。\footcites[56, 59--61]{kelsen2015}[65, 68--70]{kelsen1929}

\subsection{Schmitt 的核心批評}

Schmitt 反對把這套法階層說成從憲法一路到判決的「一貫階層序列」（\emph{durchgängige Stufenfolge}）。他的理由是：憲法與法律、法律與判決之間，不只是同一種規範依序具體化；立法者依憲法作成政治決定，法官則是依據法律對具體案件作成裁判。不能因兩者都可以說是「依據」上位規範，就把它們當成同一種法技術活動。\footcites[109--111]{schmitt2015}[38--40 n.~2]{schmitt1931}

Schmitt 尤其反對把憲法法院的活動理解成單純涵攝：當憲法律本身內容有疑義，法院若以有拘束力的方式確定它究竟包含什麼，實際上是在決定憲法內容；「任何機關只要以有權方式消除法律內容的疑義，在實質上就是以立法者身分行事」。若由法院作成，甚至是憲法立法，而非司法。\footcites[116--120]{schmitt2015}[45--48]{schmitt1931}

所以 Schmitt 說 Kelsen 的形式法學會把真正的政治決定藏在「規範位階」「法的適用」與「司法形式」後面，形成一種司法形式的虛假規範性。這不是說 Schmitt 否認所有法律拘束，而是他認為在憲法內容高度開放或存在政治衝突時，最後問題是誰有權作成決定，不是法院能否假裝以純粹法律涵攝得出答案。

\subsection{兩人的真正分歧}

兩人不是單純爭論「法院是不是法院」，而是在爭論憲法爭議的核心：

\begin{itemize}
\tightlist
\item
  Kelsen：憲法法院以實證憲法為基準，撤銷違憲一般規範；其撤銷效果具有一般性，但功能主要是受憲法拘束的法適用。
\item
  Schmitt：若法院必須先決定開放的憲法內容或政治權限邊界，這種決定就不能僅靠司法形式掩飾；它可能已是政治性、甚至憲法立法性的決定。
\end{itemize}

因此，「形式法學」在本週脈絡中最好理解成 Kelsen 對法的創設形式、位階、程序與可歸屬機關的理論；Schmitt 的批評則指出，這種形式模型在憲法真正發生政治衝突時，不能自動消除「誰來決定」的問題。

\qaentry{2026-10-12}{憲法法院守護憲法、獨立性與「立法性」}

\begin{qaprompt}
Kelsen 以憲法法院守護憲法，確保一般規範不抵觸憲法；憲法法院的獨立性；憲法法院的工作本質是立法，有一般性
\end{qaprompt}

\begin{qaquestion}
Kelsen 如何把憲法法院的憲法保障功能、機關獨立性與其撤銷一般規範所具有的立法性連在一起？
\end{qaquestion}

\subsection{守護憲法的功能}

Kelsen 的核心命題是：憲法若只能宣示立法程序與法律內容的要求，卻沒有撤銷違憲行為的制度，從法技術上就不具有完全拘束力。最有效的憲法保障，是由憲法法院撤銷違憲的法律或其他法行為，使一般規範不能在法秩序中繼續抵觸憲法。\footcites[43--45]{kelsen2015}[51--53]{kelsen1929}

因此，憲法法院不是抽象地「守護憲法精神」，而是以具體的撤銷權保障憲法對下位法規範的效力：憲法是審查基準，違憲的一般規範是審查客體，撤銷判決是使憲法拘束力實現的法律效果。

\subsection{為什麼必須獨立}

若由立法者自己決定是否撤銷自己制定的法律，憲法對立法者的拘束就會落空。因此，撤銷違憲行為必須由一個與作成違憲行為的機關完全不同，且獨立於該機關的機關直接作成。這是 Kelsen 所說有效合憲性保障的必要條件。\footcites[42]{kelsen2015}[50]{kelsen1929}

這裡的獨立性不是指憲法法院不受任何法律拘束；恰好相反，憲法法院必須高度受憲法拘束。獨立的是它與立法者及其他公權力之間的組織與決定關係，使它能對外部機關的規範作成撤銷判決，而不由受審查者自行控制審查結果。

\subsection{「本質是立法」要改成什麼}

不能簡化寫成「憲法法院的工作本質是立法」。Kelsen 的精確區分是：

\begin{itemize}
\tightlist
\item
  撤銷一般規範具有立法性：立法是創設一般規範；憲法法院撤銷法律也以一般性方式影響不特定多數案件，所以「廢止法律本身就是立法功能」，憲法法院因而是消極立法者。\footcites[46]{kelsen2015}[54]{kelsen1929}
\item
  憲法法院的功能主要仍是法的適用：國會在憲法框架內有較大的政策形成空間；憲法法院則主要適用憲法，受憲法實質拘束。因此 Kelsen 說，消極立法者「其功能在本質上由憲法決定」，功能與一般法院相近，是真正的裁判。\footcites[47]{kelsen2015}[56]{kelsen1929}
\end{itemize}

所以最準確的總結是：憲法法院撤銷一般規範的效果具有一般性，因而具有功能上的立法性；但其活動的本質拘束來自憲法，主要仍是以憲法為基準的裁判，而不是自由創設政策的積極立法。

\qaentry{2026-10-12}{憲法法院的司法性質與政治性爭議；腳註四的改引}

\begin{qaprompt}
司法性質的，不是具有政治性的

腳註四，老師說，老師會改成美國憲法條文。
\end{qaprompt}

\begin{qaquestion}
這裡應如何修正憲法法院的性質？老師所說腳註四改引美國憲法條文，應如何記錄？
\end{qaquestion}

\subsection{Kelsen 的立場：憲法法院首先是司法性質的裁判機關}

不能把「憲法法院具有一般規範效力」直接寫成「憲法法院是政治機關」。Kelsen 的制度論證是：憲法法院雖然撤銷法律的結果具有一般性，並且因此帶有消極立法者的性質，但其功能主要是憲法規範的適用，所以與一般法院相近，屬於真正的裁判。\footcites[47]{kelsen2015}[56]{kelsen1929}

因此，前一則的 Gesetzeskraft 是裁判的法律效果，不是把裁判本身改造成政治決定。應分開：

\begin{itemize}
\tightlist
\item
  性質：憲法法院是以司法程序適用憲法的裁判機關；
\item
  效果：其裁判在法律明文規定的範圍內，對一般規範及其他國家機關發生拘束或法律效力；
\item
  爭議：Schmitt 認為，若案件涉及憲法內容的疑義，而法院不是適用既有、明確的法律，而是自行決定其內容，這個決定就超出司法，實質上具有政治或立法性質。\footcites[101]{schmitt2015}[31--32]{schmitt1931}
\end{itemize}

所以課堂上的精確說法是：Kelsen 主張憲法法院的制度功能具有司法性質；Schmitt 則質疑，對憲法內容疑義作有拘束力的最終決定，能否仍稱為司法。 不是把 Kelsen 的憲法法院直接改寫成「政治性的法院」。

\subsection{腳註四}

腳註四依老師課堂指示，記為：老師將改以美國憲法的具體條文作為引據。 在老師尚未指定條文前，不先擅填條號；正式改稿時再以老師指定的美國憲法條文替換原腳註。

\qaentry{2026-10-12}{德國基本法第94條第4項與《聯邦憲法法院法》第31條}

\begin{qaquestion}
德國《聯邦憲法法院法》（BVerfGG）哪一條規定憲法法院裁判的拘束力或法律效力？這是否就是記憶中的「憲法法院判決具有法的效力」？
\end{qaquestion}

\subsection{GG Art. 94 Abs. 4 與 BVerfGG §31 的原文}

% 原稿此處列本機檔 1132_German_Legal/課堂筆記/04_修法94_n.F./Art_93-94_Source.md，排版本刪去路徑。
原文是《德國基本法》Art. 94 Abs. 4；BVerfGG §31 再將其中的拘束效力與法律效力具體化：

\begin{quote}
\emph{Die Entscheidungen des Bundesverfassungsgerichts binden die Verfassungsorgane des Bundes und der Länder sowie alle Gerichte und Behörden. Ein Bundesgesetz bestimmt, in welchen Fällen seine Entscheidungen Gesetzeskraft haben.}
\end{quote}

本機課堂翻譯也明確寫成：「聯邦憲法法院之判決對聯邦與各邦憲政機關，以及所有法院與行政機關，具有拘束力。聯邦法律另定於何種情形下，該等判決具有法律效力。」

這裡的 Gesetzeskraft 是你記得的「法律效力」：BVerfGG §31 Abs. 1 是 Bindungswirkung（拘束效力），§31 Abs. 2 才具體規定在列舉程序中裁判具有 Gesetzeskraft。它不是 Rechtskraft（既判力），也不是普通法院判決可能產生的事實上先例影響，而是法律明文賦予憲法法院裁判的特殊一般效力。

所以要分三層：

\begin{itemize}
\tightlist
\item
  Rechtskraft：個案裁判的既判力；
\item
  Bindungswirkung（§31 Abs. 1）：對憲法機關、法院及行政機關的拘束；
\item
  Gesetzeskraft（§31 Abs. 2）：法律在特定憲法訴訟程序中賦予裁判的法律效力，尤其使違憲法律的判斷不只是該案中的不適用。
\end{itemize}

這也正好接回 Schmitt 所談的普通法院與憲法法院之差異：普通法院在具體案件中不適用普通法律，至多形成事實上的先例影響；憲法法院依第31條第2項取得的 Gesetzeskraft，才是制度上明文規定的、超出單一案件的裁判效力。

因此，這一題的完整定位是：1132 德文法學 → GG Art. 94 Abs. 4 → BVerfGG §31 → Gesetzeskraft。

\qaentry{2026-10-12}{個案拘束力、先例效力與權威性解釋}

\begin{qaprompt}
個案拘束力？權威性的解釋？
\end{qaprompt}

\begin{qaquestion}
一般法院不適用違憲法律時，個案拘束力、先例或權威性解釋，以及一般性憲法裁判效力，應如何區分？
\end{qaquestion}

\subsection{三種不同層次}

\begin{itemize}
\tightlist
\item
  個案拘束力：法院在具體案件中，對當事人及該案件的裁判結果發生拘束；普通法律在該案中不予適用，但法律本身沒有因此被一般性消滅。
\item
  先例／權威性解釋的效力：上級法院或最高法院的理由與法律見解，可能對其他法院形成先例上的拘束、說服或「或多或少可計算的影響」；這是裁判理由的外溢效力，不等於法律本身失效。
\item
  一般性規範失效：憲法法院以集中審查判決撤銷違憲一般規範，效力直接指向規範本身，通常及於不特定案件；這才是 Kelsen 所說憲法法院作為消極立法者的核心效果。
\end{itemize}

Schmitt 對帝國法院1925年判決的說明，正是前兩者而非第三者：帝國法院在具體案件中優先適用憲法律，普通法律只在該案中不適用；該裁判作為判例，對其他法院有一定但非規範一般失效的影響。\footcites[83, 88--89]{schmitt2015}[16, 20]{schmitt1931}

\subsection{理論上的個案效力與事實上的一般拘束力}

必須再區分兩個角度：

\begin{itemize}
\tightlist
\item
  理論上／形式上：一般法院判決的直接效力只及於該案，普通法律仍留在法秩序中；即使最高法院拒絕適用，也不是正式宣告法律對所有人失效。
\item
  事實上／實際上：最高法院判決會透過先例、法院層級、下級法院的遵循與法律人的預期，對其他案件產生普遍而穩定的影響，甚至形成近似一般拘束的效果。
\end{itemize}

因此，「事實上的一般拘束力」不是「理論上的一般規範效力」。前者是判決理由與司法權威的外溢，後者才是憲法法院以判決直接改變一般規範效力範圍的制度效果。Schmitt 說帝國法院判決作為先例，對其他法院只有「或多或少可計算的影響」，正是在描述前者。

\subsection{「權威性解釋」的特殊用法}

若說的是 1831 年薩克森憲法第153條的 \emph{authentische Interpretation}（權威性解釋），意思是：政府與等級會議無法協議消除憲法文書疑義時，由國事法院作成並要求遵行的解釋。Schmitt 指出，任何機關只要以有權方式消除法律內容的疑義，在實質上就是以立法者身分行事；若確定的是憲法律內容，甚至是憲法立法，而不是單純司法。\footcites[117]{schmitt2015}[45]{schmitt1931}

所以「權威性解釋」不能只看名稱：若只是個案中提供判決理由，屬先例效力；若對憲法內容作一般且有拘束力的確定，可能已超出普通司法，具有立法或憲法立法性。

\qaentry{2026-10-12}{美國司法審查的「英美法系土壤」與老師文章}

\begin{qaprompt}
英美法系的土壤？／老師要寫的文章\ldots\ldots 等《權力分立新論》改版，這學期應該能出來
\end{qaprompt}

\begin{qaquestion}
Schmitt 所說美國最高法院特殊地位的「英美法系土壤」是什麼？老師提到的文章與《權力分立新論》改版，應如何記錄？
\end{qaquestion}

\subsection{「英美法系土壤」不是地理比喻}

這裡的「土壤」是指一整套制度與法律文化條件，而不是單純指英國、美國位於歐洲以外。Schmitt 認為，美國最高法院的地位是在盎格魯撒克遜式的司法國家（\emph{Justizstaat}）中形成；其特徵包括：

\begin{itemize}
\tightlist
\item
  普通法院對公共生活具有廣泛控制；
\item
  司法審查附隨於真實的 \emph{cases or controversies}，而非抽象地回答政治機關的諮詢；
\item
  普通法與判例傳統使法院能在具體爭訟中逐步形成憲法判斷；
\item
  美國沒有歐陸式行政法國家所形成的同樣機關分工與法院結構。
\end{itemize}

因此，美國最高法院的憲法審查是一種在特定法律土壤中長成的制度。把它抽離出來，直接移植到有行政法、行政國家、不同憲法結構與政治衝突的歐陸國家，制度名稱可以相同，實際功能卻未必相同。\footcites[80--82]{schmitt2015}[12--14]{schmitt1931}

\subsection{老師文章與改版資訊}

課堂上提到：老師另有文章準備撰寫，並表示等待《權力分立新論》改版；依課堂說法，本學期應有機會出版。這裡只能記為老師在課堂上提供的出版預期，不能先寫成已出版、已定稿或已確定的書目資料。正式引用仍須等實際出版資訊確定後，再依文章或新版書籍的正式書名、刊物、日期與頁碼處理。

\qaentry{2026-10-12}{Schmitt：美國最高法院模式不能全然移植到歐陸}

\begin{qaprompt}
Schmitt：\ldots\ldots 不能全然移植到歐洲大陸的可能性
\end{qaprompt}

\begin{qaquestion}
Schmitt 為什麼認為美國最高法院的憲法審查模式，不能直接移植到歐洲大陸國家？
\end{qaquestion}

Schmitt 的原文不是抽象地說「美國制度比較好、歐洲制度比較差」，而是說美國最高法院的特殊地位，不能脫離其歷史與制度條件移植：

\begin{quote}
„eine Interpretation, die nicht einfach auf die sozial und politisch ganz anders geartete Lage eines kontinental-europäischen Staates übertragen werden kann``

「這種解釋不能直接移植到社會與政治狀況完全不同的歐陸國家。」
\end{quote}

他提出的差異至少有四層：

\begin{itemize}
\tightlist
\item
  國家類型不同：美國最高法院發展於盎格魯撒克遜式的司法國家（\emph{Justizstaat}）；歐陸國家具有行政法與行政國家的傳統。
\item
  法院功能不同：美國最高法院原則上只裁判真實的 \emph{cases or controversies}，不是德國意義上的國事法院（\emph{Staatsgerichtshof}），也不提供抽象諮詢意見。
\item
  審查基準不同：Schmitt 認為美國法院運用財產、價值、自由等一般原則，保護一套原則上沒有爭議的社會與經濟秩序；這不能直接套入威瑪德國的政治與社會衝突。
\item
  憲法概念不同：美國脈絡下的「憲法」主要可由自由市民法治國的基本權、人身自由與私有財產理解；歐陸的憲法還涉及行政國家、政黨國家、聯邦權限與整體政治秩序。
\end{itemize}

所以 Schmitt 的「不能全然移植」不是說歐洲完全不能有司法審查，而是說：同一個制度名稱，放在不同國家類型、法院結構、行政法傳統與憲法內容中，功能可能完全不同。把美國最高法院直接當作歐陸憲法法院的模型，正是他所反對的「不加檢驗的移植與神話化」。\footcites[80--82]{schmitt2015}[12--14]{schmitt1931}

\qaentry{2026-10-12}{Schmitt 對美國的理解，不等於老師對美國的完整理解}

\begin{qaprompt}
師：Schmitt 對美國的理解和老師對美國的理解不完全一樣／師：不僅限於經濟領域
\end{qaprompt}

\begin{qaquestion}
Schmitt 如何理解美國最高法院的憲法審查？為什麼不能把他的描述理解成美國憲法只保護經濟領域？
\end{qaquestion}

Schmitt 的原文是在作一個特定的制度比較：美國最高法院發展於盎格魯撒克遜式的司法國家（\emph{Justizstaat}），普通法院對公共生活具有廣泛控制；在他的概括中，美國最高法院依第五、第十四修正案及 \emph{due process of law}，保護自由市民法治國所理解的基本權，尤其是人身自由與私有財產，並對抗立法者，維護市民的社會與經濟秩序。\footcites[80--82]{schmitt2015}[13--16]{schmitt1931}

但這段話有三個限制：

\begin{itemize}
\tightlist
\item
  它是 Schmitt 為了對比德國帝國法院與美國最高法院所作的歷史---制度類型化，不是美國憲法審查全部內容的清單。
\item
  Schmitt 說的是「社會與經濟秩序」（\emph{soziale und wirtschaftliche Ordnung}），不是只說經濟利益；其前文已包括人身自由、私有財產與基本權的自由主義法治國秩序。
\item
  老師提醒的重點是：理解美國不能被 Schmitt 的自由市民法治國框架限縮。美國憲法審查還涉及平等、政治參與、言論與宗教自由、權力分立以及其他非經濟性的基本權問題；不能把 Schmitt 的選擇性描述當成美國憲法的完整圖像。
\end{itemize}

所以應分開兩句：Schmitt 描述的是他用來反駁「德國國事法院＝美國最高法院」的美國類型；老師談的則是較完整、較不受十九世紀自由市民法治國框架限制的美國憲法審查。「不僅限於經濟領域」不是否定 Schmitt 原文，而是防止把他的比較法概括誤讀成美國憲法審查的全部範圍。

\qaentry{2026-10-12}{1925年11月4日帝國法院判決（RGZ 111, 322）}

\begin{qaprompt}
1925年判決
\end{qaprompt}

\begin{qaquestion}
Schmitt 所說的 1925 年11月4日帝國法院判決，確立了什麼？它是否使帝國法院成為憲法守護者？
\end{qaquestion}

這是德國帝國法院（\emph{Reichsgericht}）1925年11月4日的判決，刊於 RGZ 111, 322。判決承認：法官受帝國法律拘束（威瑪憲法第102條），並不排除法官在普通法律與具有優先效力的憲法律衝突時，拒絕適用普通法律。

判決的德文核心表述是：

\begin{quote}
\emph{Die richterliche Bindung an das Gesetz schließt nicht aus, daß einem Reichsgesetz oder einzelnen seiner Bestimmungen vom Richter die Gültigkeit insoweit aberkannt werden kann, als es mit anderen, den Richter ebenfalls bindenden Vorschriften im Widerspruch steht.}
\end{quote}

意思是：法官受法律拘束，並不排除法官在普通法律與其他同樣拘束法官的規範牴觸時，在該限度內否認其效力。\footcites[82]{schmitt2015}[15]{schmitt1931}

但 Schmitt 對這個判決作了限縮解讀：

\begin{itemize}
\tightlist
\item
  法官只能在具體案件中，把事實涵攝於相互衝突的法律規定，優先適用憲法律；
\item
  這是普通法律在該案中的拒絕適用／不予適用（\emph{Nichtanwendung des einfachen Gesetzes auf den konkreten Fall}）；
\item
  拒絕適用不等於否定法律效力：普通法律仍留在法秩序中，只是在本案不作為裁判基礎；
\item
  不是宣告普通法律一般性、對世地失效；
\item
  這種審查是附隨於個案裁判的、分散的司法審查，原則上任何法官都可能作成。
\end{itemize}

Schmitt 因此說，不能因帝國法院承認這種審查權，就稱它為嚴格意義上的憲法守護者。法院仍然受法律拘束，也沒有因此取得一般性的憲法守護或政治決斷權。\footcites[82--84, 88--90]{schmitt2015}[12, 16, 20--22]{schmitt1931}這正是「一般法院」與 Kelsen 式集中型憲法法院之間不能直接類比的關鍵。

\qaentry{2026-10-12}{一般法院與憲法法院：不要犯類比錯誤}

\begin{qaprompt}
一般的法院？「很容易掉入一個類比的錯誤」
\end{qaprompt}

\begin{qaquestion}
Schmitt 為什麼說一般法院與憲法法院不能直接類比？一般法院在個案中不適用違憲法律，是否就等於憲法法院撤銷法律？
\end{qaquestion}

原文首先就把這個誤解寫出來：

\begin{quote}
„Die prozeßentscheidenden Gerichte der bürgerlichen, strafrechtlichen oder verwaltungsrechtlichen Gerichtsbarkeit sind nicht in einem präzisen Sinne Hüter der Verfassung. Der Irrtum, sie als solche zu bezeichnen, liegt allerdings nahe \ldots``

「民事、刑事或行政審判中作成程序終局決定的法院，並非嚴格意義上的憲法守護者。把它們稱為憲法守護者的錯誤，確實很容易發生\ldots\ldots」
\end{quote}

Schmitt 接著說，原因是一般法院可能行使所謂的實質司法審查權：審查普通法律是否符合憲法律，衝突時不適用普通法律。\footcites[12]{schmitt1931}[82]{schmitt2015}

不能等同。Schmitt 所說的一般法院，是在具體案件中依憲法律審查普通法律；若兩者衝突，法院在該案中選擇適用憲法律，將普通法律不予適用（\emph{Nichtanwendung des Gesetzes}）。這嚴格說不是宣告普通法律自始或一般性失效，而是個案裁判的法律效果。\footcites[82--83, 88--89]{schmitt2015}[15--16, 20]{schmitt1931}

憲法法院則是集中審查一般規範，並以判決對規範本身作一般性的撤銷；這是 Kelsen 所說的憲法法院作為消極立法者。兩者都會碰到「法律是否符合憲法」，但審查方式、標的、效力與制度位置不同。

因此，類比錯誤在於：看見一般法院也能不適用違憲法律，就推論一般法院已經等同憲法法院；或看見憲法法院也是法院，就推論它只是在個案中作普通司法。Schmitt 正是在阻止這種跳躍：司法獨立與個案不適用，是一般法院有限的防禦性審查，不等於全面的憲法守護或一般規範撤銷權。

這也要和 Kelsen 分開：Kelsen 將兩者放在同一法創設階層中，強調功能上的連續性；Schmitt 則強調制度與功能的實質差別，避免以「都是法院」「都在適用憲法」作過度類比。

\qaentry{2026-10-12}{司法獨立不等於憲法守護者}

\begin{qaprompt}
Schmitt 認為：司法獨立／司法可以守護憲法的一部分，不等於憲法守護者
\end{qaprompt}

\begin{qaquestion}
為什麼法院具有司法獨立，或能守護關於司法自身的憲法規範，仍不等於它就是一般性的憲法守護者？
\end{qaquestion}

Schmitt 的區分是：法院可以是憲法某一部分的守護者，但不因此取得整體憲法的守護者地位。法院基於司法獨立，可以防禦其他國家權力對法院自身基礎、組織地位與審判權的侵犯；這是對司法自身憲法地位的防禦性保護。

但「憲法守護者」在制度意義上，必須能在憲政衝突中代替一般人民的不服從與抵抗，並對整體政治秩序作具有代表性的決定。法院的獨立性只說明它在審判時不受其他權力指揮，不能自動推出它有權裁判所有政治性憲法爭議。Schmitt 因此說，法院即使能守護自身所依據的那部分憲法，也不足以構成一般性的 \emph{Hüter der Verfassung}（憲法守護者）。\footcites[85, 89--90]{schmitt2015}[17, 21--22]{schmitt1931}

最精確的表述是：司法獨立是司法守護憲法的條件之一，但不是成為全面憲法守護者的充分條件。

\qaentry{2026-10-12}{司法作為憲法的守護者：Schmitt 的限定命題}

\begin{qaprompt}
司法作為憲法的守護者
\end{qaprompt}

\begin{qaquestion}
Schmitt 是否完全否認司法可以作為憲法的守護者？
\end{qaquestion}

不是完全否認，而是作嚴格限縮。Schmitt 認為，法院可以守護憲法中構成司法自身基礎與地位的部分，尤其是司法獨立、法院組織地位及防止其他國家權力侵犯司法權的規範。這種守護是法院為維持自身憲法地位所作的防禦性審查；不能由此推論法院就是一般性的憲法守護者。\footcites[85]{schmitt2015}[17]{schmitt1931}

所以要區分兩個命題：

\begin{itemize}
\tightlist
\item
  可以：法院拒絕適用直接侵害司法獨立或法院憲法地位的法律，守護憲法中關於司法自身的部分。
\item
  不可以直接推出：法院因此有權審查並決定所有政治性憲法爭議，成為全面的憲法守護者。
\end{itemize}

Schmitt 的重點不是「司法完全不能守護憲法」，而是司法守護的客體與範圍有限；一旦爭議涉及整體政治秩序、權力配置或憲法內容的根本決定，就不能只以司法獨立或法院形式正當化法院的全面介入。

\qaentry{2026-10-12}{Schmitt 並不以憲法法院為憲法守護者的理想中心}

\begin{qaprompt}
Carl Schmitt 沒有那麼著迷於憲法法院
\end{qaprompt}

\begin{qaquestion}
Schmitt 是否和 Kelsen 一樣，把憲法法院視為憲法守護的核心機關？
\end{qaquestion}

不是。Schmitt 的主要工作不是替憲法法院設計更好的審查制度，而是質疑：憲法守護者的功能，能否一般性地以司法形式行使。他認為，若法院面對的是憲法內容本身的政治疑義、權力衝突或高度開放的原則，法院不能只靠獨立性與司法程序，就把政治決定轉化成司法裁判。\footcites[95, 100--103]{schmitt2015}[26, 31--33]{schmitt1931}

Schmitt 特別反對把法院塑造成能裁判一切憲法疑義的機關。這樣的國事法院將成為與國會、帝國總統及帝國政府並列的政治機關；問題不是它是否使用法院名稱，而是它實際上取得了政治決定權。\footcites[100--101]{schmitt2015}[31]{schmitt1931}

\subsection{Schmitt 的替代方案}

在威瑪憲法的脈絡中，Schmitt 主張由帝國總統擔任憲法守護者，理由是帝國總統具有全民選舉的代表性、政治中立的定位，以及在憲政危機中作成維持整體政治統一之決定的能力。這不是 Kelsen 式的集中型憲法法院，而是政治性、代表性的憲法守護模式。

因此，本週兩人的對立不能整理成「兩人都重視憲法法院，只是對法院功能的理解不同」。較準確是：

\begin{itemize}
\tightlist
\item
  Kelsen：以獨立的憲法法院集中審查並撤銷違憲一般規範，讓憲法具有完整拘束力。
\item
  Schmitt：懷疑法院能否承擔一般性的憲法守護；法院至多在真正法律性、可由明確規範裁判的範圍內活動，憲法整體的守護則涉及政治決斷與代表性。
\end{itemize}

所以 Schmitt 並非「著迷於憲法法院」，而是在批判把憲法守護法院化、司法化與規範技術化；他真正關心的是憲法危機中誰能作最後的政治決定。

\qaentry{2026-10-12}{Schmitt 文中的「帝國憲法」}

\begin{qaprompt}
帝國憲法？
\end{qaprompt}

\begin{qaquestion}
Schmitt 文中所說的「帝國憲法」（\emph{Reichsverfassung}）是哪一部憲法？
\end{qaquestion}

這裡指的是1919 年《德意志國憲法》（\emph{Verfassung des Deutschen Reichs}），通常稱威瑪憲法（\emph{Weimarer Reichsverfassung}），是威瑪共和國的憲法，不是 1871 年德意志帝國成立時的憲法。

「\emph{Reich}」在德文法律與政治語境中是當時德國國家名稱的一部分，故 \emph{Reichsverfassung} 直譯為「帝國憲法」；放到威瑪時代，較不會誤解的中譯是威瑪憲法／德意志國憲法。Schmitt 討論憲法守護者時，所涉及的是威瑪憲法第 48 條總統緊急權、第 25 條解散國會、第 54 條不信任案，以及國會、帝國總統與帝國政府之間的權限配置。

因此，前述「帝國總統」也是威瑪憲法下的德意志國總統（\emph{Reichspräsident}），不是君主，也不是 1871 年帝國皇帝。Schmitt 主張由這個具有全國代表性的共和國元首擔任憲法守護者，正是針對威瑪憲政危機提出的制度論證。

\qaentry{2026-10-12}{Schmitt 的刑事責任：政治犯罪、部長責任與憲法守護}

\begin{qaprompt}
Schmitt 刑事責任
\end{qaprompt}

\begin{qaquestion}
Schmitt 如何看政治犯罪的刑事審判與部長刑事責任？它們能否成為憲法守護的方式？
\end{qaquestion}

Schmitt 把刑事責任放在司法守護憲法的最低限度、事後性範圍內。叛國罪等政治犯罪的刑事審判，可以在廣義上算是憲法保護，因為它制裁已經完成的危害國家行為；但它只能處理「已經終結、屬於過去的事實」，不能事前處理政治秩序或正在發生的憲法衝突。\footcites[95--96]{schmitt2015}[26--27]{schmitt1931}

部長法律責任也有同樣問題。十九世紀君主立憲時期，部長刑事責任原本被期待用來控制政府；但 Schmitt 認為，與部長的政治責任相比，司法責任實際作用很小。當政治責任被改寫成刑事司法問題，司法形式反而成為使政治責任失去作用的可靠方式：

\begin{quote}
„Die Justizförmigkeit war das sichere Mittel der politischen Wirkungslosigkeit.``

「司法形式是確保政治上不起作用的可靠手段。」
\end{quote}

因此，Schmitt 不是說刑事責任沒有意義，而是說它不能被抽象地當作「法治國」的充分答案。部長彈劾若涉及內容不明、政治權限分配或憲法本身的疑義，法院不能只因取得審判權，就自行補充憲法內容；否則刑事審判會轉化為政治決定，形成司法的政治化。\footcites[90--93, 96--99]{schmitt2015}[22--29]{schmitt1931}

所以 Schmitt 的結論是：刑事責任可以事後制裁明確的政治犯罪，但不等於一般性的憲法守護者，也不能以刑事法院取代對政治性憲法衝突的制度決定。

\qaentry{2026-10-12}{Schmitt 第107頁：憲法律有疑義時，司法的界限}

\begin{qaprompt}
p.107 憲法法律產生疑義時\ldots\ldots{}
\end{qaprompt}

\begin{qaquestion}
Schmitt 在第107頁說憲法律產生疑義時，為什麼仍要堅持法律與判決的實質區別？
\end{qaquestion}

這裡的正確用語是憲法律（\emph{Verfassungsgesetz}），不是泛稱「憲法法」。Schmitt 在本節開頭提出兩個互相連結的命題：

\begin{quote}
„keine unabhängige Justiz ohne inhaltliche Bindung an ein Gesetz, keine inhaltliche Bindung an das Gesetz ohne sachliche Verschiedenheit von Gesetz und Richterspruch gibt``

「沒有在內容上受法律拘束，就沒有獨立的司法；法律與判決若沒有實質區別，就沒有內容上的法律拘束。」
\end{quote}

意思是：法官之所以是獨立的法官，不是因為他可以自由決定一切，而是因為他受一個內容上可確定的法律拘束。法律與判決若沒有實質區別，法院就不再是依法律裁判，而是在自行創設原本應由立法者決定的內容。

因此，即使憲法律產生疑義，也不能只因法院獨立，就把政治性的內容決定交給法院。法官可以在法律容許的範圍內進行解釋，但不能取得屬於立法者的政治決定權；否則就會改變司法在憲法上的地位。Schmitt 因而主張，立法、執行與裁判之間仍須有實質區分；若完全沒有這種區分，國家就不具有自由主義意義的憲法。\footcites[36]{schmitt1931}[107--108 n.~55]{schmitt2015}

這也是他批評 Kelsen「從憲法到判決的一貫階層序列」的切入點：Kelsen 把憲法、法律、判決放進同一個法創設階層；Schmitt 則說，憲法律若本身內容有疑義，法院有拘束力地決定其內容，實質上就是立法，甚至是憲法立法，而不再只是司法。\footcites[109, 116--120]{schmitt2015}[38, 45--48]{schmitt1931}

\qaentry{2026-10-12}{腳註71的三種妥協與 Schmitt 的多重假定}

\begin{qaprompt}
三篇文章做了三種不同妥協；很多的假定
\end{qaprompt}

\begin{qaquestion}
腳註71所引 Kelsen 三部作品中的「妥協」是否是同一種妥協？Schmitt 把它們連在一起時，預設了哪些假定？
\end{qaquestion}

不能把三部作品簡化成 Kelsen 反覆主張同一個「妥協」。較準確是：妥協在三個文本中具有不同的理論功能。

\begin{enumerate}
\def\labelenumi{\arabic{enumi}.}
\tightlist
\item
  國家理論的妥協（《一般國家學》）：議會國家中的國家意志形成，表現為多數與少數、不同政治力量之間的協商與妥協；妥協是議會制度的運作方式。
\item
  民主理論的妥協（《民主的本質與價值》）：民主不被理解為毫無拘束的多數統治，而是多數與少數持續妥協；妥協是民主正當性的實質條件。
\item
  憲法裁判的妥協（〈國事裁判的本質與發展〉）：憲法裁判保障憲法、保護少數，並使多數與少數之間的持續妥協不被多數專政破壞；妥協是憲法裁判的政治功能。
\end{enumerate}

Schmitt 在腳註71所做的，不是證明三部作品使用完全相同的「妥協」，而是把三種功能串成一個批判對象。要完成這個串接，他至少作了幾個假定：

\begin{itemize}
\tightlist
\item
  議會的妥協性可以代表民主的本質；
\item
  民主的多數／少數妥協可以由憲法裁判制度保存；
\item
  憲法裁判保護少數，等於保護多元政治力量之間的妥協；
\item
  Kelsen 對妥協的政治、民主與裁判功能，都可用同一套規範理論說明；
\item
  憲法、憲法律、規範、政治決定與契約／妥協之間的差別，可以由「規範」概念連接起來。
\end{itemize}

Schmitt 的批判力正來自這裡：他揭示 Kelsen 理論之間的連續性，但也可能把不同文本、不同層次的「妥協」壓縮成一條過於整齊的推論鏈。故腳註71既是文獻串聯，也是大量理論前提的集中處；讀它時必須分開文本中 Kelsen 實際說了什麼與Schmitt 如何把三處材料組合成批判。

\qaentry{2026-10-12}{老師說「腳註71非常精彩」：Schmitt 如何借 Kelsen 攻擊規範主義}

\begin{qaprompt}
師：腳註71非常精彩
\end{qaprompt}

\begin{qaquestion}
Schmitt 的腳註71為什麼精彩？它在憲法妥協、多元主義與 Kelsen 的憲法裁判理論之間做了什麼連結？
\end{qaquestion}

腳註71 的精彩之處，在於 Schmitt 沒有只抽象批評 Kelsen，而是直接把 Kelsen 三部作品放在一起讀：

\begin{itemize}
\tightlist
\item
  《一般國家學》（\emph{Allgemeine Staatslehre}）第324頁；
\item
  《民主的本質與價值》（\emph{Wesen und Wert der Demokratie}）第57頁；
\item
  〈國事裁判的本質與發展〉（\emph{Wesen und Entwicklung der Staatsgerichtsbarkeit}）第81頁。
\end{itemize}

Schmitt 由此指出，Kelsen 把現代民主理解為多數與少數之間的持續妥協，又把聯邦國的憲法裁判與這種民主妥協連結起來。換言之，在 Kelsen 那裡，憲法裁判不只是審查規範，也具有保護少數、維持多元政治妥協的功能。\footcites[123]{schmitt2015}[72]{kelsen2015}[81]{kelsen1929}

Schmitt 的反擊是：Kelsen 的報告把憲法＝憲法律＝規範，於是把原本不同的三件事------政治決定、法律規範、契約／妥協------都用 \emph{Norm} 這個詞包起來。這會遮蔽憲法究竟是誰作成的政治決定、哪些集團只是談判妥協，以及法院日後是否正在替憲法補上原本沒有決定的內容。

因此，腳註71 不是普通的文獻註，而是 Schmitt 對 Kelsen 的理論系譜追擊：從民主理論的「妥協」、國家理論的「議會妥協」，一路追到憲法裁判的少數保護，藉此主張 Kelsen 的規範主義把決斷與政治衝突隱藏在規範概念之後。它正好把本節的核心問題集中起來：威瑪憲法究竟是人民的政治決定，還是多元團體之間可由法院繼續協商、解釋與補充的契約／妥協文件。

\qaentry{2026-10-12}{Schmitt：憲法契約與憲法妥協}

\begin{qaprompt}
憲法契約／憲法妥協
\end{qaprompt}

\begin{qaquestion}
Schmitt 所批評的「憲法契約」與「憲法妥協」有何不同？為什麼他認為這種理解不適合威瑪憲法？
\end{qaquestion}

\subsection{\texorpdfstring{憲法契約（\emph{Verfassungsvertrag}）}{憲法契約（Verfassungsvertrag）}}

憲法契約觀把國家看成若干具有自己權利的當事人所組成：例如君主與人民、國王與議會兩院、政府與代議機關之間，以契約共同決定憲法。依此理解，憲法爭議就是當事人之間的權利爭執，修憲是重新締結契約，憲法疑義則先由當事人協議解決，協議不成才交由法院。\footcites[120--122]{schmitt2015}[40--43]{schmitt1931}

\subsection{\texorpdfstring{憲法妥協（\emph{Verfassungskompromiss}）}{憲法妥協（Verfassungskompromiss）}}

憲法妥協觀更進一步把憲法看成多元政治、宗教、階級與利益團體之間的談判結果：憲法本身以及國家意志形成，都只是多元勢力之間的妥協。這裡的妥協有兩種可能：

\begin{itemize}
\tightlist
\item
  真正的妥協：各方對某一共同決定達成實質協議。
\item
  形式妥協／延宕妥協（\emph{dilatorische Formelkompromisse}）：各方無法作成真正決定，便用模糊、可供相反立場共同接受的公式，把衝突延後，留待日後由政治機關或法院決定。
\end{itemize}

Schmitt 對威瑪憲法第二部的批評，主要針對第二種：許多條文沒有真正決定，只是把衝突包裝成憲法文字；日後若由法院確定其內容，法院就會第一次創設原本不存在的規範，實質上成為憲法立法者。\footcites[119--120]{schmitt2015}[47--48]{schmitt1931}

\subsection{Schmitt 的結論}

Schmitt 認為，威瑪憲法的正確理解不是各政治集團彼此締結的契約，也不是各利益團體共同持有的妥協文件；它是德意志人民作為同質且不可分割的政治統一體，依制憲權以積極政治決定為自己制定的憲法。\footcites[122]{schmitt2015}[62]{schmitt1931}

所以，Schmitt 的問題不是反對一切政治協商，而是反對把憲法的效力基礎降格為「各方只要曾經談成，就各自擁有一部分憲法」。憲法是共同政治統一體的決定；妥協若只是延後決定，不能取代這項決定。

\qaentry{2026-10-12}{那麼 Schmitt 認為誰是憲法守護者？}

\begin{qaprompt}
那麼什麼是呢？
\end{qaprompt}

\begin{qaquestion}
既然一般法院、帝國法院與單純遵守憲法的機關或人民，都不是嚴格意義上的憲法守護者，那麼 Schmitt 認為誰才是？
\end{qaquestion}

在 Schmitt 的威瑪憲法方案中，真正的憲法守護者是帝國總統（\emph{Reichspräsident}），不是一般法院，也不是一個可以裁判所有政治疑義的憲法法院。帝國總統之所以具有守護者地位，不是因為他能像法院一樣涵攝法律，而是因為他具有制度上的代表性、相對中立的地位，以及在憲政衝突或危機中維持整體政治統一的決斷能力。

「憲法守護者」的制度意義，在 Schmitt 那裡是要取代人民各自的不服從與抵抗：必須有一個特定機關，能在整體憲政秩序受到威脅時作成具有拘束力的政治決定。這就是他主張由帝國總統承擔守護功能的理由。\footcites[22ff.]{schmitt1931}[90, ch.~6]{schmitt2015}

因此，本週兩人的答案要明確分開：

\begin{itemize}
\tightlist
\item
  Kelsen：憲法守護者是獨立的憲法法院，以撤銷違憲一般規範保障憲法。
\item
  Schmitt：憲法守護者是帝國總統，以代表性與政治決斷維持整體憲政秩序；法院只能在明確法律問題與自身司法地位的防禦範圍內活動。
\end{itemize}

Schmitt 的「憲法守護者」不是最會解釋憲法的法院，而是能在政治統一體面臨衝突時，代表整體作成最後決定的制度性機關。

\qaentry{2026-10-12}{Schmitt 認為誰不是憲法守護者？}

\begin{qaprompt}
Schmitt 認為什麼不是憲法守護者
\end{qaprompt}

\begin{qaquestion}
Schmitt 認為哪些機關不能因為適用憲法或審查法律，就被稱為嚴格意義上的憲法守護者？
\end{qaquestion}

首先是民事、刑事、行政的一般法院，包括帝國法院（\emph{Reichsgericht}）。它們可以在具體案件中審查普通法律與憲法律是否衝突，並在衝突時不適用普通法律；但這只是附隨於個案裁判的司法審查，不是對一般規範作一般性撤銷。因此，Schmitt 說它們「不是在精確意義上的憲法守護者」。\footcites[12]{schmitt1931}[82--83]{schmitt2015}

其次，單純遵守憲法的一般原則，也不會形成制度意義上的憲法守護者。如果只要遵守憲法的國家機關或人民都算守護者，那麼每一個機關、甚至每一個公民都會變成潛在的守護者；這最多導出一般的不服從權或抵抗權，沒有建立一個特定的制度機關。\footcites[89--90]{schmitt2015}[21--22]{schmitt1931}

所以 Schmitt 的判準是：個案不適用普通法律，不等於一般性守護憲法；司法獨立，也不等於全面的憲法守護者地位。

\qaentry{2026-10-12}{\emph{Reasonable disagreement}：合理分歧不是任何分歧}

\begin{qaprompt}
\emph{reasonable disagreement}
\end{qaprompt}

\begin{qaquestion}
\emph{reasonable disagreement} 是什麼？它如何區分釋字第 748 號與黨產條例前案？
\end{qaquestion}

\subsection{意義}

\emph{Reasonable disagreement} 是「合理之人之間的分歧」，不是「只要有人不同意，就都算合理」。它承認憲法問題常有多種可由法律理由支持的答案；因此，法院判決即使可能是錯的，也不當然等於逾越權限或欠缺法律正當性。Fallon 的界線是：判決是否仍在 the bounds of reasonable legal judgment 之內。\footcite[38, 40]{kelsen2015}

判斷「合理」至少要看兩面：

\begin{itemize}
\tightlist
\item
  言之有理：論證有法律理由可循，不是只宣示結論或以政治立場代替法律論證。
\item
  言之有物：有憲法文字、法律規範、判例、制度脈絡或案件事實作為實質依據。
\end{itemize}

所以，合理分歧不等於多數意見必然正確，而是表示少數意見不能僅因敗訴就把多數意見視為濫權；反過來，若判決以牽強說理、無視明確先例或沒有實質法律依據，就可能逸出合理法律判斷的範圍。

\subsection{748 號與黨產條例前案}

課堂上把釋字第 748 號歸為 reasonable disagreement：即使對其憲法解釋與制度效果可以有不同評價，多數意見仍有憲法權利、平等保障及法律論證作為依據，屬於可爭辯但仍在合理法律判斷範圍內的分歧。

相對地，老師把黨產條例前案的不受理決議稱為 unreasonable disagreement，重點不是「意見不同」本身，而是認為多數大法官以牽強說理，輕率廢棄長期形成的解釋先例，因而不是正常的憲法解釋分歧，而是欠缺足夠法律理由的偏離。\footcite[38, 40]{fallon2018}

\subsection{與 Kelsen 的連結}

這個概念可與 Kelsen 限制憲法法院審查基準的主張並讀：憲法法院應以實證法為基準，不能把法官多數自己認為的正義，直接當成撤銷國會法律的憲法規範。Kelsen 指出，法院多數對「正義」的看法可能與人民多數及國會多數完全相反；因此憲法若要授權法院審查法律內容，原則、指針與界限應盡可能明確。\footcites[59--61]{kelsen2015}[68--70]{kelsen1929}

\qaentry{2026-10-12}{美國的 constitutional amendment：有沒有不能修的內容？}

\begin{qaprompt}
美國的amendment
\end{qaprompt}

\begin{qaquestion}
美國憲法的 \emph{amendment} 是否也有不可修改的憲法內容？
\end{qaquestion}

美國《憲法》第5條的明文限制，比德國《基本法》第79條第3項窄得多：任何州未同意，不得剝奪其在參議院的平等投票權（equal Suffrage in the Senate）。這是美國憲法明文保留的不可修改事項。第5條另有一個只適用於1808年以前的暫時限制，現在已經失效。

所以三國的問題要分開：

\begin{itemize}
\tightlist
\item
  美國：有一項明文、具體的程序／聯邦結構保障------未經該州同意，不得取消其參議院平等代表權；沒有像德國 Art. 79 Abs. 3 那樣列出民主、基本權與法治的廣泛永恆條款。
\item
  德國：Art. 79 Abs. 3 明文保護聯邦結構、各邦參與立法，以及 Art. 1、20 的基本原則。
\item
  我國：沒有美國第5條或德國第79條第3項那種完整明文；釋字第499號是由憲法基本原則與正當修憲程序發展出修憲界限。
\end{itemize}

因此，美國的 \emph{amendment} 不是「什麼都能修」，但其明文不可修的範圍主要是州在參議院的平等代表權，而不是一整套由法院概括宣告的基本憲法核心。美國是否存在超出第5條文字的「基本結構」限制，學理上有爭論，不能直接當成美國現行法的確定答案。\footcites{conan-artV1}{conan-artV5}

\qaentry{2026-10-12}{美國 constitutional amendments 的歷次修改與違憲審查}

\begin{qaprompt}
美國的 constitutional amendment（憲法修正）的歷次修改？對美國的 constitutional amendment（憲法修正）的違憲審查？
\end{qaprompt}

\subsection{一、歷次修正的整體輪廓}

美國自1789年憲法生效後，國會共提出33件修正案，州只批准其中27件；所以現在的美國憲法有 27件正式修正案。第5條的州議會修正案會議提案途徑至今沒有使用過；實際上全部修正案都是由國會兩院提出。

依歷史階段整理如下：

\begin{itemize}
\tightlist
\item
  第1至第10修正（1791）：權利法案（Bill of Rights），保障宗教、言論、出版、集會、刑事程序、財產等自由。
\item
  第11修正（1795）：限制以州為被告的聯邦訴訟；第12修正（1804）：調整總統與副總統的選舉程序。
\item
  第13至第15修正（1865--1870）：重建修正案，廢除奴隸制度、保障平等與正當法律程序、禁止以種族等理由剝奪男性投票權。
\item
  第16至第19修正（1913--1920）：所得稅、參議員直選、禁酒、婦女投票權。
\item
  第20至第22修正（1933--1951）：總統任期開始、總統任期限制等政府運作規則；其中第21修正撤銷第18修正的全國禁酒。
\item
  第23至第27修正（1961--1992）：哥倫比亞特區總統選舉人票、總統繼任與失能、投票年齡、國會議員薪資生效時點等。
\end{itemize}

這裡最重要的例子是：第21修正以同樣的憲法修正程序撤銷第18修正。所以美國的修正案不是永遠不可變；後來的修正案可以改變或撤銷先前的修正案，但仍受第5條明文限制拘束。

\subsection{二、美國最高法院是否審查「修正案本身」？}

要把「違憲審查」拆成三種：

\begin{enumerate}
\def\labelenumi{\arabic{enumi}.}
\tightlist
\item
  審查修正案的提出、批准程序：可以審查，但最高法院對國會與州批准程序保留很大的政治部門空間。
\item
  審查修正案的實質內容：第5條明文限制以外，最高法院沒有建立像德國 Art. 79 Abs. 3 那樣的廣泛永恆條款審查；也沒有宣告任何已完成的聯邦修正案因實質違憲而無效。
\item
  審查依修正案制定的普通法律或行政行為：可以審查。修正案成為憲法的一部分後，執行它的普通法律仍須符合其他憲法規範。
\end{enumerate}

\subsection{三、幾個關鍵判例}

\begin{itemize}
\tightlist
\item
  National Prohibition Cases, 253 U.S. 350 (1920)：最高法院確認第18修正案已成為憲法的一部分，並認為禁止酒精飲料的內容在第5條所保留的修憲權範圍內。這是「法院沒有把修正案當成普通法律來審查其違憲」的最重要例子。
\item
  Dillon v. Gloss, 256 U.S. 368 (1921)：承認國會得為修正案的批准設定合理且確定的期間。
\item
  United States v. Sprague, 282 U.S. 716 (1931)：批准方式由國會在州議會批准與州批准會議之間選擇。
\item
  Coleman v. Miller, 307 U.S. 433 (1939)：州議會批准程序中的若干問題屬政治性問題，法院不取代國會作最後判斷。
\end{itemize}

因此，美國的模式不是「最高法院完全不能碰修正案」，而是：法院可以處理第5條程序性問題，但對修正案內容的審查極為有限；已完成的修正案原則上具有憲法效力，法院主要審查其實施法律，而不是把修正案本身當成普通法律宣告違憲。

\subsection{四、和德國及我國的對照}

\begin{itemize}
\tightlist
\item
  美國：第5條只有具體的明文不可修改事項（未經該州同意，不得剝奪其參議院平等代表權）；第18修正可由第21修正撤銷。
\item
  德國：Art. 79 Abs. 3 明文保護聯邦結構、各邦立法參與，以及 Art. 1、20 的基本原則，形成實質性的 Ewigkeitsgarantie。
\item
  我國：釋字第499號直接把修憲程序及憲法基本原則納入審查，並宣告第五次修憲的部分規定無效；這比美國最高法院目前對修正案本身所採的態度更積極。
\end{itemize}

所以在本週的 Kelsen 問題上，美國提供的是一個重要對照：美國有憲法至上與司法審查，但不當然等於法院有權審查修正案的一切內容；修正權的界限首先由 Article V 的文字決定。\footcites{conan-artV1}{conan-artV5}{conan-amdt18-6}

\qaentry{2026-10-12}{哪些憲法內容不能修改？1132 德文法學提到的 Art. 79 Abs. 3}

\begin{qaprompt}
憲法上哪些東西是不能修的？看1132法學德文有講到嗎？我有印象
\end{qaprompt}

\begin{qaquestion}
憲法修正權有沒有不可逾越的界限？1132 的法學德文課程是否講過這件事？
\end{qaquestion}

\subsection{1132 本機筆記的確有這個脈絡}

你的印象是對的。1132「05\_文章／文章.md」在 Verfassungskern vs. Gesetzesdetail 後面接著寫：

\begin{itemize}
\tightlist
\item
  Verfassungsänderung（憲法修正）與 Ewigkeitsgarantie（永恆條款）；
\item
  本段雖未直接處理《德國基本法》第79條第3項，但涉及法院制度，仍須符合第20條原則。
\end{itemize}

也就是說，課堂筆記當時是在用「憲法核心／法律細節」區分：制度的基本結構應由憲法規定；可變動的程序、配額等細節則可交給普通法律。這不是把所有憲法文字都凍結，而是辨認哪些基本原則不能被修憲掏空。

來源：\texttt{1132/1132\_German\_Legal/課堂筆記/05\_文章/文章.md}。

\subsection{德國基本法的明文答案：Art. 79 Abs. 3 GG}

《德國基本法》第79條第3項（Ewigkeitsklausel／Ewigkeitsgarantie）規定，修正基本法而涉及下列事項者，不許可：

\begin{enumerate}
\def\labelenumi{\arabic{enumi}.}
\tightlist
\item
  聯邦劃分為各邦的基本結構；
\item
  各邦在聯邦立法中的基本參與；
\item
  《基本法》第1條及第20條所規定的原則。
\end{enumerate}

其中第1條的核心包括人的尊嚴不可侵犯，以及基本權對立法、行政、司法的直接拘束；第20條的核心包括民主、社會國、聯邦國、人民主權、權力分立與法治拘束等基本秩序。第79條第3項的文字是「涉及」（\emph{berührt werden}），所以重點是不得動搖這些原則的核心，不是說第1條、第20條中的每一個制度細節都完全不能調整。\footcite[Art. 79]{gg}

\subsection{和本週 Kelsen、Schmitt 及釋字第499號的關係}

這正好接到本週的問題：如果修憲法律本身也是法律，誰來審查它？

\begin{itemize}
\tightlist
\item
  德國基本法直接設置第79條第3項，承認修憲權本身受憲法核心拘束；這是明文的「不可修」界限。
\item
  Kelsen 的形式階層模型本身，若沒有更高位階的明文規範，較難單獨推出一個實質上不可修改的憲法核心。這就是前面所說的理論漏洞：修憲法律若形式上有效，究竟依何規範判斷其逾越修憲權？
\item
  我國憲法沒有完全對應德國第79條第3項的明文永恆條款；釋字第499號因此從憲法整體基本原則及正當修憲程序，建立對修憲權的限制。這是我國大法官的憲法解釋，不是直接把德國條文搬過來。
\end{itemize}

因此，這一題要分清楚：德國是「明文的 Ewigkeitsgarantie」；我國釋字第499號則是由憲法解釋發展出修憲界限；Kelsen 的問題在於，形式法學若沒有更高規範，不能自行產生這個實質核心。

\qaentry{2026-10-12}{修憲法律也是法律，但什麼不能審查？}

\begin{qaprompt}
修改憲法的法律也是法律／什麼不能審查
\end{qaprompt}

\begin{qaquestion}
修憲法律在形式上也是法律，為什麼 1925 年帝國法院判決仍把它排除在審查範圍之外？
\end{qaquestion}

\subsection{形式上的法律，不等於同一種審查客體}

修憲法律（\emph{Verfassungsgesetz} 或 \emph{verfassungsänderndes Gesetz}）當然也是由法律程序制定的規範；但在 1925 年帝國法院判決所表明的審查射程中，法院的審查權只及於帝國普通法律，不及於依威瑪憲法第76條程序制定的修憲法律。這首先是該判決對自身審查權的限制，不是可以脫離脈絡、直接當成所有憲法法院的普遍命題。

因此要分開三件事：

\begin{itemize}
\tightlist
\item
  形式分類：修憲法律仍是法律；
\item
  位階與功能：它是用來修改憲法的憲法律，不是普通法律；
\item
  審查射程：普通法院不能因認為修憲法律內容不妥，就對其作一般性違憲審查。依 Schmitt 的整理，修憲權是否有界限、第76條程序是否被濫用，都不在帝國法院這種狹窄審查權內。\footcites[84]{schmitt2015}[16]{schmitt1931}
\end{itemize}

所以「修憲法律也是法律」只能推出它具有法律規範性，不能推出普通法院當然可以審查它。在這個歷史脈絡下，不能審查的是：依憲法修憲程序制定的修憲法律本身，尤其其內容是否逾越修憲權界限。

\subsection{見解的歸屬}

\begin{itemize}
\tightlist
\item
  直接來源：1925 年 11 月 4 日帝國法院判決（RGZ 111, 322）對普通法院審查權的自我限縮；
\item
  本文中的位置：Schmitt 在第一章轉述這個判決，並以它說明德國帝國法院的審查只是附隨、分散的個案審查，不能與美國最高法院的審查或 Kelsen 式集中型憲法法院類比；\footcite[82--84, 86--90]{schmitt2015}
\item
  不是本文的無條件普遍命題：Schmitt 並非在此抽象宣稱任何憲法法院永遠只能審查普通法律，而是在描述帝國法院判決的歷史射程，並把它放入自己反對「一般法院即憲法守護者」的論證；
\item
  Kelsen 的不同問題：Kelsen 主張由集中型憲法法院審查一般規範，審查客體與普通法院的附隨審查不能直接混同。
\end{itemize}

\subsection{不要和 Kelsen 的國際法問題混在一起}

Kelsen 的討論另保留一個不同層次：如果採國際法秩序優位，修憲法律也可能因牴觸條約或一般國際法而有瑕疵；但那是以國際法優位及憲法承認國際法為前提的問題，不是 1925 年帝國法院依威瑪憲法第76條所採的普通法院審查射程。\footcites[57--59]{kelsen2015}[66--68]{kelsen1929}

\qaentry{2026-10-12}{第四次修憲「凸槌、擦槍走火」：從常設國大走向任務型國大}

\begin{qaprompt}
師：第四次修憲是凸槌的，是擦槍走火的。非常危險，

常設國大改成
任務型國大
\end{qaprompt}

\begin{qaquestion}
老師為什麼把第四次修憲評為「凸槌」「擦槍走火」「非常危險」？常設國民大會改成任務型國民大會，是第四次修憲的結果嗎？
\end{qaquestion}

\subsection{先校正制度時間線}

老師這句話是在評價一條憲政改革路徑，但制度上要分開：

\begin{itemize}
\tightlist
\item
  第四次修憲（1997）：仍維持常設性國民大會，主要處理精省、行政院長同意權、倒閣與總統解散立法院等中央政府體制問題；同時重新安排國民大會代表的產生方式與任期。
\item
  第五次修憲（1999）：進一步設計依立法委員得票比例產生國民大會代表，並牽涉國代、立委任期延長；這部分因程序及實質憲法衝突被釋字499宣告失效。
\item
  第六次修憲（2000）：才正式把國民大會改造成任務型國民大會：國大不再是常設機關，只有在立法院提出修憲、領土變更或總統副總統彈劾案時，才依比例代表方式選出代表，於特定任務完成後終止。
\item
  第七次修憲（2005）：廢除國民大會，改由人民以公民投票複決修憲案與領土變更案。
\end{itemize}

因此，「常設國大改成任務型國大」若問直接的法律變更，答案是第六次修憲；若問老師為什麼把它放在第四次修憲的危險後果脈絡中，答案是：第四次修憲先大幅改造中央政府與國大制度，留下了後續修憲機關自我改造、擴張或縮減權限的高度風險。

\subsection{「擦槍走火」的憲法意義}

這裡的危險不只是修正某個機關名稱，而是常設的修憲機關同時成為憲政權力競爭者：它可以透過修憲改變自己的組成、職權與任期，甚至可能把原本的人民授權轉化為政黨或現任代表的自我延續。第五次修憲的國代延任案，正把這個危險具體引爆，最後由釋字499踩下煞車。

所以可把老師的話整理成：第四次修憲不是單純「成功建立任務型國大」；它是一次制度設計失控的警訊。後來第六次修憲以任務型國大收束常設國大的權力，第七次修憲再把國民大會整體廢除。

\subsection{這份 PDF 在材料中的內容位置}

這篇文章不是李登輝或蘇南成的回憶錄，而是事後的政治史／憲法學分析。它的核心問題是：第五、六次修憲如何圍繞國民大會的國會改革展開，以及政治權力的推動邏輯如何和修憲的法律正當性發生衝突。

\begin{itemize}
\tightlist
\item
  頁111--113：提出全文的基本問題------修憲權掌握在政治權力競逐者手中時，政治邏輯如何轉化為符合立憲主義的法律；這是理解「凸槌、擦槍走火」的理論入口。
\item
  頁120--121：進入第五次修憲的政黨比例代表問題，包括未參加政黨者能否取得國大代表資格，以及政黨得票能否替代國民對修憲代表的直接授權；這一段直接接到你前面說的「政黨主權／國民主權」。
\item
  頁165--166：回到修憲程序理性與制度後果，說明第五次修憲為何引發釋字499，以及第六次修憲如何以任務型國大處理常設國大的制度困局。
\end{itemize}

所以，這份 PDF 最適合放在本題的分析層：它解釋為什麼第四次修憲留下危險的制度路徑，以及第五次、六次如何爆發與收束。若要找「蘇南成親自怎麼說李登輝為什麼要修憲」，則應回到國科會成果報告裡的〈蘇南成先生訪問錄〉；該報告現已有本機掃描檔可供核對。

\subsection{國科會成果報告中的〈蘇南成先生訪問錄〉：李登輝為何推動修憲}

《第五、六次修憲之研究》（NSC 89-2414-H-002-032）所收〈蘇南成先生訪問錄〉，於2000年11月15日在臺南進行，訪問者為李炳南、曾建元、顏韶儀。訪談位於刊本頁 293--303。\footcite[頁 293--303]{su2000interview}

與「李登輝為什麼修憲」直接相關的原文可分為四個層次：

\begin{itemize}
\tightlist
\item
  蘇南成說，第五次修憲一方面是要把前四個階段零散的增修條文作一次系統整理，形式上可以思考為「擴大增修」、特別法或基本法；另一方面是要把選舉次數「化繁為簡」。當時延任構想原本同時涉及總統、立法委員與國民大會。針對外界認為延任案是李登輝要延長任期，蘇南成明確表示：「有的人認為這是李登輝要出來延任，其實李登輝是絕對不願意的，但你不願意也要接受。」\footcite[頁 295--296]{su2000interview}
\item
  蘇南成指出，李登輝直到1999年8月19日才明確表示總統不能延任。其說法將時間順序與政治背景連在一起：1999年7月9日「兩國論」引發兩岸緊張，美國政府介入後，若再推動延任，容易使外界理解為「為了總統要延任」，因而增加政治與國際上的疑慮。\footcite[頁 297]{su2000interview}
\item
  蘇南成回憶，1999年8月18日連戰幕僚擬妥延任案並交由民進黨提出；民進黨擔心李登輝不同意，遂派人到臺南請許文龍向李登輝報告。蘇南成於8月19日上午見李登輝，下午再見連戰。這段是對修憲案政治運作與李登輝本人態度的當事人敘述。\footcite[頁 300]{su2000interview}
\item
  對於李登輝是否參與「臺灣基本法」構想，蘇南成先答「都是傳聞，沒有啦」，隨後補充：「李總統本身也希望力求政治改革，在他有生之年，在他任期之內，能為台灣制定一套更完整的增修條文，這是他的一個心願。」這是訪談中最直接把李登輝的修憲動機表述為「完成較完整的增修條文」的一段。\footcite[頁 301]{su2000interview}
\end{itemize}

因此，這份訪談所呈現的不是單一的「李登輝想延任」解釋，而是三層動機交疊：整理前階段增修條文的制度性動機、簡化選舉與任期安排的制度動機，以及在任期內完成較完整憲政改革的政治歷史想像。至於延任案，蘇南成的說法是李登輝原本不願意，後來因政治情勢而必須面對；這是當事人訪談中的敘述，仍須與其他訪談及同時期文件交叉比對，不能直接當作無爭議的客觀定論。

你找到的 PDF 是李炳南、曾建元，〈第五、六次修憲國會制度改革相關議題之政治邏輯與法理分析〉，《人文及社會科學集刊》第14卷第1期（2002），頁111--173；它是刊登版研究論文，不是李炳南主持的完整國科會成果報告，也不是〈蘇南成先生訪問錄〉本身。完整報告為李炳南主持之國科會專題研究計畫成果報告《第五、六次修憲之研究》（NSC 89-2414-H-002-032，2001）。文章摘要明確把第五次修憲經釋字499失效，以及第六次改成任務型國大，連成同一條制度脈絡。\footcites{lee2002}{pres325}{pres93}{jy499}

\qaentry{2026-10-12}{人民不入黨的自由、政黨主權與國民主權：釋字499踩下煞車}

\begin{qaprompt}
人民有不參加政黨的自由，怎麼反應意見？政黨主權？國民主權？憲法被顛覆了？

大法官是對的，進來，踩剎車
\end{qaprompt}

\begin{qaquestion}
若國民大會代表改由政黨得票比例「附帶產生」，會不會把國民主權轉成政黨主權？釋字第499號是否正是在這裡介入制止？
\end{qaquestion}

\subsection{不是「反對政黨」，而是反對政黨取代人民的直接授權}

人民有組織政黨的自由，也有不加入政黨、以獨立身分表達政治意見與參與公共事務的自由。政黨可以整合意見、提出政綱、參與選舉，但政黨的得票不能自動等於人民已直接選出另一個具有修憲權的機關。

第五次修憲的問題在於：國民大會代表本身沒有另行舉行選舉，卻要依立法委員選舉中政黨及獨立參選人的得票比例分配名額；而國民大會又具有修改憲法、複決修憲案及其他重大憲政職權。這會使人民對立法委員的投票，被轉換成對另一個修憲機關的間接投票，形成政黨或政黨得票替代人民直接授權的效果。

\subsection{499 的法律語言：國民主權與自由民主憲政秩序}

釋字第499號不是用一句「政黨主權違憲」作結，而是從憲法基本原則推出限制：

\begin{itemize}
\tightlist
\item
  憲法第1條的民主共和國原則；
\item
  憲法第2條的國民主權原則；
\item
  第二章人民權利的保障；
\item
  權力分立與制衡；
\item
  由上述原則形成的自由民主憲政秩序。
\end{itemize}

大法官明確說，這些原則是憲法規範秩序存立的基礎；憲法設置的機關也受其拘束。若修憲使未經選舉、由政黨得票附帶產生的人員行使本質上必須由民意代表行使的修憲與政權職權，就不是單純調整組織，而是與民主憲政秩序發生規範衝突。

大法官對獨立參選人的說明尤其關鍵：獨立參選人不屬於任何黨派或政團，沒有共同政策綱領；若依附他人得票取得國大代表資格，候選人便無從以自己的理念與主張訴諸選民而獲選，這不符合憲法保障人民參政權的意旨（釋字499理由第10段）。

\subsection{「踩剎車」的真正意思}

所以你說「大法官是對的，進來，踩剎車」，法律上可整理成：當修憲機關把政黨得票轉成未經獨立選舉的國民大會代表，並讓其繼續行使修憲等重大權力時，大法官以憲法基本原則制止修憲權對國民主權的替代。

「憲法被顛覆」是對制度效果的強烈概括；較精確的法學表述是：第五次修憲相關設計已使修憲機關的民主正當性、人民直接授權與憲法基本原則之間產生不可容許的規範衝突。釋字499所說的「憲法自我防衛」，就是在這個意義下的憲法內部煞車機制。\footcite[理由書第 9 至 11 段]{jy499}

\qaentry{2026-10-12}{國民大會的「政黨代表」：如果今天還存在，會聽到誰？}

\begin{qaprompt}
國民大會的「政黨代表」，聽賴清德 黃國昌 鄭麗文
\end{qaprompt}

\begin{qaquestion}
老師說「不然今天還看到國民大會」，是否應具體理解成：國民大會若延續至今，會成為政黨代表直接進行憲法政治的場所？
\end{qaquestion}

對。你補的不是一般性的「政黨政治」，而是第五次修憲所設計的國民大會政黨比例代表：第四屆國民大會代表300人，依立法委員選舉中各政黨推薦候選人的得票比例分配名額；同一設計也把獨立參選人的得票放進比例分配。這個制度設計正是釋字第499號審查的對象之一。

所以你說「聽賴清德、黃國昌、鄭麗文」，是在做一個當代化的想像：如果國民大會今天仍是修憲機關，公共政治中就可能出現一個由各政黨代表組成、直接處理憲法修改與國家制度問題的場域。這裡不是說這三人曾經或必然會成為國民大會代表，而是用今天可辨認的政黨領導／代表人物，具體想像那種機關會如何被政黨政治占據。

\subsection{499 對「政黨代表」設計的精確問題}

釋字第499號指出：政黨推薦的候選人具有共同政綱與政策，可以理解為以政黨得票比例分配名額；但獨立參選人既不屬於任何黨派或政團，沒有共同政策綱領，不能只是依附他人得票而取得代表資格。這顯示第五次修憲並非單純調整名額，而是在重新安排誰能進入修憲機關、以何種政治資格代表人民。

因此，老師那句「不然今天還看到國民大會」的具體制度含義是：如果499沒有阻止第五次修憲延長國民大會代表任期，今天可能仍有一個以政黨比例代表為主要組成方式的憲法機關；我們會看到政黨代表在其中爭論修憲，而不是只有立法院內的政黨競爭。\footcites{gazette6297}{jy499}

\qaentry{2026-10-12}{老師評釋字第499號：如果沒有這個解釋，今天可能還看得到國民大會}

\begin{qaprompt}
師：499是個偉大的解釋，大法官有勇氣，至少做了對的事情

不然今天還看到國民大會
\end{qaprompt}

\begin{qaquestion}
老師為什麼把釋字第499號評為「偉大的解釋」？「不然今天還看到國民大會」指的是什麼？
\end{qaquestion}

\subsection{老師這句話的評價對象}

「偉大」「有勇氣」「做了對的事情」是老師對釋字第499號的評價，不是本文作者或整理者應擅自改寫成中性的定論。它所肯定的是：大法官在面對國民大會以修憲形式延長自身政治生命、並涉及修憲程序與憲法基本原則的案件時，沒有把「修憲法律也是憲法」當成完全不可審查，而是承認修憲權仍受憲法拘束。

\subsection{「不然今天還看到國民大會」的制度脈絡}

釋字第499號宣告第五次修憲中關於國民大會代表任期等相關規定，自解釋公布日起失其效力，並使先前的增修條文繼續適用。老師的意思是：如果大法官不承認修憲程序與修憲內容可以受到憲法審查，國民大會就可能透過延長自身任期的修憲，繼續作為實際運作中的憲法機關存在。

這是對制度後果的反事實式概括，不應改寫成「釋字499單獨造成國民大會後來全部消失」：國民大會制度後來還經歷第六次、第七次修憲的制度變更。準確的說法是，釋字499先阻止了第五次修憲以重大程序瑕疵延長國民大會代表任期，並確立修憲權不是無限制的權力。

\subsection{和本週理論問題的連接}

這正是 Kelsen 式問題在我國憲法實務中的具體化：若修憲法律形式上也是法律，仍須追問修憲機關是否依正當程序行使權限、是否已逾越憲法所容許的修改界限。釋字499的歷史重要性，不只在於處理一次國民大會任期爭議，也在於它以憲法解釋承認修憲權本身可以受到憲法控制。
